% =============================================================================
% Simula Training Data Framework — Implementation Specification
% Main document - includes chapter files
% =============================================================================
\documentclass[11pt,a4paper]{report}

\usepackage[utf8]{inputenc}
\usepackage[T1]{fontenc}
\usepackage{lmodern}
\usepackage{amsmath}
\usepackage[margin=2.5cm]{geometry}
\usepackage{graphicx}
\usepackage{booktabs}
\usepackage{longtable}
\usepackage{tabularx}
\usepackage{multirow}
\usepackage{array}
\usepackage{float}
\usepackage{enumitem}
\usepackage{xcolor}
\usepackage{tcolorbox}
\usepackage{tikz}
\usepackage[table]{colortbl}
\usetikzlibrary{arrows.meta,positioning,shapes.geometric}

\newtcolorbox{adrbox}[2][]{colback=blue!5,colframe=blue!75!black,fonttitle=\bfseries,title=ADR: #2,#1}

% --- Semantic Hyperlinking --------------------------------------------------
\usepackage{xr-hyper}
\usepackage[hidelinks]{hyperref}

% External references to peer specifications
\externaldocument[reg:]{../regulations/regulations-spec}
\externaldocument[tb:]{../tb/tb-review-spec}
\externaldocument[ar:]{../arabic/arabic-ap-spec}

\usepackage[capitalize,noabbrev]{cleveref}
\usepackage{fancyhdr}
\usepackage{titlesec}

% ─── Code listings ───────────────────────────────────────────────────────────
\usepackage{listings}

\definecolor{codebg}{RGB}{248,248,248}
\definecolor{codegreen}{RGB}{0,128,0}
\definecolor{codegray}{RGB}{128,128,128}
\definecolor{codeblue}{RGB}{0,0,180}
\definecolor{codepurple}{RGB}{128,0,128}
\definecolor{sapblue}{HTML}{0B5FA5}
\definecolor{sapgreen}{HTML}{00695C}
\definecolor{sapamber}{HTML}{B8360F}
\definecolor{sapgrey}{HTML}{546E7A}
\definecolor{sapmidnight}{HTML}{1C2B39}
\definecolor{sapgold}{HTML}{F0AB00}
\definecolor{simred}{HTML}{B71C1C}
\definecolor{simblue}{HTML}{1565C0}
\definecolor{simgreen}{HTML}{2E7D32}
\definecolor{simgrey}{HTML}{546E7A}
\definecolor{simamber}{HTML}{F0AB00}
\definecolor{hanagreen}{HTML}{2E7D32}
\definecolor{llmpurple}{HTML}{6A1B9A}
\definecolor{saplight}{HTML}{F5F5F5}

\lstdefinelanguage{Python}{
keywords={class,def,return,if,elif,else,for,in,import,from,async,await,
True,False,None,self,try,except,finally,with,as,yield,not,and,or,
Optional,Dict,List,Tuple,Any,str,int,float,bool,dataclass,field,Enum},
keywordstyle=\color{codeblue}\bfseries,
ndkeywords={@dataclass,@property,asyncio,httpx},
ndkeywordstyle=\color{codepurple}\bfseries,
identifierstyle=\color{black},
sensitive=false,
comment=[l]{\#},
commentstyle=\color{codegreen}\ttfamily,
stringstyle=\color{sapblue}\ttfamily,
morestring=[b]',
morestring=[b]"
}

\lstdefinelanguage{jsonx}{
basicstyle=\ttfamily\small,
columns=fullflexible,
breaklines=true,
showstringspaces=false,
comment=[l]{//},
morestring=[b]",
stringstyle=\color{sapgreen},
keywordstyle=\color{sapblue},
}

\lstset{
language=Python,
backgroundcolor=\color{codebg},
basicstyle=\ttfamily\small,
breakatwhitespace=false,
breaklines=true,
captionpos=b,
keepspaces=true,
showspaces=false,
showstringspaces=false,
showtabs=false,
tabsize=2,
frame=single,
rulecolor=\color{black!10}
}

% --- Glossary Helper Commands -----------------------------------------------
\definecolor{pathcolor}{HTML}{004D40}
\newcommand{\schemapath}[1]{\texttt{\textcolor{pathcolor}{#1}}}
\newcommand{\enumval}[1]{\texttt{\textcolor{sapgreen}{#1}}}
\newcommand{\fieldref}[1]{\texttt{\textcolor{sapblue}{#1}}}

\begin{document}

\title{\Huge Simula Training Data Framework \\
\textbf{Implementation Specification} \\
\large Synthetic Data Generation for Text-to-SQL \\
\small Ref: \texttt{SIM-TRAIN-2026-v1.2}}
\author{SAP-OSS AI Training Team}
\date{April 2026}

\maketitle

% --- Incorporated Executive Summary ---

\maketitle
\thispagestyle{empty}

% =============================================================================
% ABSTRACT
% =============================================================================
\begin{abstract}
\noindent
This specification defines the full implementation of the \textbf{Simula} reasoning-driven
synthetic training data generation framework for SAP HANA Cloud. The system addresses a
critical gap identified in the existing training pipeline: \emph{there is no direct pipe
to HANA to generate training datasets for model training}. Three orthogonal components
are specified in full:

\begin{itemize}[nosep]
\item a \textbf{HANA Schema Extraction Pipeline} that connects to SAP HANA Cloud via
the AI Core PAL MCP service, queries \texttt{SYS.TABLE\_COLUMNS} for live metadata,
and populates a schema registry;
\item a \textbf{Taxonomy Generation Engine} that implements Algorithm~1 from the Simula
research paper---factor identification, breadth-first expansion with Best-of-N
proposals, and generator-critic refinement;
\item a \textbf{Training Data Generation Engine} that implements Algorithm~2---taxonomic
mix sampling, meta-prompt generation, complexification, and double-critic filtering.
\end{itemize}

\noindent
All three components consume a single vLLM TurboQuant backend via the OpenAI-compatible API
and output JSONL training data suitable for LoRA fine-tuning. The specification targets
three application domains: Arabic document processing (invoice extraction), Trial Balance
review (variance analysis and LLM commentaries), and ESG analytics (financed emissions
classification).

\medskip
\noindent\textbf{Reading order.} Chapter~\ref{chap:overview} situates the work;
Chapter~\ref{chap:schema} defines the data structures; Chapter~\ref{chap:extraction}
covers the HANA extraction pipeline; Chapter~\ref{chap:taxonomy} covers the taxonomy
engine; Chapter~\ref{chap:generator} covers the data generator; Chapter~\ref{chap:config}
specifies environment and CLI; Chapter~\ref{chap:evaluation} covers the evaluation
framework; Chapter~\ref{chap:references} provides references and traceability.

\medskip
\noindent\textbf{Conventions.} File paths are typeset in \texttt{monospace} and are
relative to the repository root \texttt{/Users/user/Documents/sap-oss}. Enum values
in running text appear in \textsc{small caps}. Python code listings include line numbers
for cross-reference.
\end{abstract}

\tableofcontents
\listoftables
\listoffigures
\newpage

% =============================================================================
% EXECUTIVE SUMMARY
% =============================================================================
\section*{Executive Summary}
\addcontentsline{toc}{section}{Executive Summary}

\subsection*{Business Problem}

The SAP OSS Studio AI platform requires high-quality training data to fine-tune
Text-to-SQL models for domain-specific applications (Trial Balance, Arabic invoices,
ESG analytics). The current pipeline uses \textbf{static Excel files} as schema sources,
creating three critical issues:

\begin{enumerate}[nosep]
\item \textbf{Schema drift:} Training data becomes stale when HANA schemas evolve.
\item \textbf{Coverage gaps:} Manual templates miss edge cases and rare query patterns.
\item \textbf{Quality variance:} No systematic quality assurance for generated examples.
\end{enumerate}

\subsection*{Solution: Simula Framework}

This specification implements the \textbf{Simula} framework from Davidson et al. (2026),
a research-validated methodology for reasoning-driven synthetic data generation. The
framework addresses all three issues through:

\begin{itemize}[nosep]
\item \textbf{HANA-Direct Extraction} — Live schema queries via AI Core PAL MCP service.
\item \textbf{Taxonomic Coverage} — Hierarchical factor trees ensure systematic diversity.
\item \textbf{Agentic Quality Control} — Double-critic filtering with Elo-calibrated complexity.
\end{itemize}

\subsection*{Key Metrics and Targets}

\begin{table}[H]
\centering
\begin{tabular}{llll}
\toprule
\textbf{Metric} & \textbf{Paper Target} & \textbf{Implementation} & \textbf{Source} \\
\midrule
Global diversity & 0.48--0.54 cosine & Target: 0.52 & Figure 4 \\
Local diversity & 0.10--0.30 cosine & Target: 0.20 & Figure 4 \\
Complexity ratio & 50\% (c=0.5) & Adaptive 0.2--0.5 & Figure 7 \\
Critic acceptance & p(y) $>$ 0.8 & Threshold: 0.5--0.7 & Figure 3 \\
Taxonomy completeness & 74\% (Simula) & Min: 70\% & Table 2 \\
Taxonomy soundness & 97\% (Simula) & Min: 90\% & Table 2 \\
Saturation ratio & 83\% & 83\% & Section 4.3 \\
\bottomrule
\end{tabular}
\end{table}

\subsection*{Expected Outcomes}

\begin{enumerate}[nosep]
\item \textbf{10$\times$ training data} with consistent quality (100k+ examples).
\item \textbf{Full taxonomic coverage} — no systematic gaps in query patterns.
\item \textbf{Calibrated complexity} — downstream model benefits from high-complexity data.
\item \textbf{Audit trail} — every example traceable to taxonomy nodes and meta-prompts.
\end{enumerate}

\newpage
% =============================================================================
% CHAPTER 1: OVERVIEW AND SCOPE
% =============================================================================
\chapter{Overview and Scope}
\label{chap:overview}

\section{The Problem}

The SAP OSS training pipeline requires high-quality Text-to-SQL training data that reflects
the actual schema structure of SAP HANA Cloud tables. The current implementation in
\texttt{src/training/pipeline/} uses metadata from \textbf{Excel files}
(\texttt{DATA\_DICTIONARY.xlsx}, \texttt{ESG\_DATA\_DICTIONARY.xlsx}) that were originally
exported from HANA or defined manually.

While the \texttt{HanaClient} in \texttt{src/training/pipeline/hana\_client.py} has the
capability to query \texttt{SYS.TABLE\_COLUMNS} for live schema metadata, this HANA-direct
extraction is \textbf{not currently integrated} into the main training data generation
pipeline.

\begin{table}[H]
\centering
\caption{Current vs. Target Training Data Pipeline}
\label{tab:gap-analysis}
\begin{tabularx}{\textwidth}{lXX}
\toprule
\textbf{Dimension} & \textbf{Current State} & \textbf{Target State} \\
\midrule
Schema Source & Static Excel/CSV files & Live HANA Cloud metadata \\
Extraction Method & Manual export & MCP-based automated query \\
Training Data & Template-based expansion & Reasoning-driven synthesis \\
Diversity Control & None & Taxonomic coverage \\
Quality Assurance & None & Double-critic filtering \\
Complexity Control & None & Calibrated Elo scoring \\
\bottomrule
\end{tabularx}
\end{table}

\section{Goals}

The specification targets three outcomes:

\begin{enumerate}[nosep]
\item \textbf{HANA-Direct Extraction.} Every training dataset is generated from live
schema metadata queried via the AI Core PAL MCP service.
\item \textbf{Taxonomic Coverage.} The conceptual space of the target domain is mapped
using hierarchical taxonomies, ensuring global diversity.
\item \textbf{Reasoning-Driven Synthesis.} Each training example is generated from
first principles using agentic refinement, ensuring quality, local diversity,
and controlled complexity.
\end{enumerate}

\section{In-Scope Domains}

\begin{table}[H]
\centering
\caption{Application Domains}
\label{tab:domains}
\begin{tabularx}{\textwidth}{llX}
\toprule
\textbf{Domain} & \textbf{HANA Schema} & \textbf{Training Task} \\
\midrule
Arabic Document Processing & \texttt{AP\_INVOICES}, \texttt{VAT\_CHECKLISTS} & Invoice field extraction, VAT compliance \\
Trial Balance Review & \texttt{TB\_VARIANCES}, \texttt{GL\_BALANCES} & Variance analysis, commentary generation \\
ESG Analytics & \texttt{ESG\_METRICS}, \texttt{FINANCED\_EMISSIONS} & Scope classification, net-zero alignment \\
\bottomrule
\end{tabularx}
\end{table}

\section{System Context}

\begin{figure}[H]
\centering
\begin{tikzpicture}[
node distance=1.2cm and 2cm,
box/.style={rectangle, draw, rounded corners, minimum width=2.5cm, minimum height=1cm, align=center, font=\small},
hana/.style={box, fill=hanagreen!20, draw=hanagreen},
mcp/.style={box, fill=sapblue!20, draw=sapblue},
llm/.style={box, fill=llmpurple!20, draw=llmpurple},
output/.style={box, fill=sapgold!20, draw=sapgold},
arr/.style={-{Stealth[length=3mm]}, thick},
]
% HANA Layer
\node[hana] (hana) {SAP HANA\\Cloud};

% MCP Layer
\node[mcp, right=of hana] (mcp) {AI Core PAL\\MCP :8084};

% Extraction
\node[box, right=of mcp, fill=saplight] (extract) {Schema\\Extractor};

% Registry
\node[box, below=of extract, fill=saplight] (registry) {Schema\\Registry};

% LLM
\node[llm, right=2.5cm of extract] (llm) {vLLM\\TurboQuant\\:9180};

% Taxonomy
\node[box, below=of llm, fill=saplight] (taxonomy) {Taxonomy\\Builder};

% Generator
\node[box, below=of taxonomy, fill=saplight] (generator) {Data\\Generator};

% Output
\node[output, right=of generator] (output) {Training Data\\(.jsonl)};

% Arrows
\draw[arr] (hana) -- node[above, font=\tiny]{SYS.TABLE\_COLUMNS} (mcp);
\draw[arr] (mcp) -- node[above, font=\tiny]{execute\_sql} (extract);
\draw[arr] (extract) -- (registry);
\draw[arr] (registry) -- (taxonomy);
\draw[arr, dashed] (llm) -- node[right, font=\tiny]{Best-of-N} (taxonomy);
\draw[arr] (taxonomy) -- (generator);
\draw[arr, dashed] (llm) -- node[right, font=\tiny]{Critic} (generator);
\draw[arr] (generator) -- (output);

% Labels
\node[below=0.3cm of hana, font=\tiny\color{hanagreen}] {Database};
\node[below=0.3cm of mcp, font=\tiny\color{sapblue}] {Protocol};
\node[below=0.3cm of llm, font=\tiny\color{llmpurple}] {LLM Backend};
\end{tikzpicture}
\caption{System context. HANA Cloud (green) provides live schema metadata via the MCP
protocol (blue). The vLLM backend (purple) powers taxonomy expansion and data generation.
Output is JSONL training data (gold).}
\label{fig:system-context}
\end{figure}

\section{Research Paper Traceability}

This specification implements the Simula framework from Davidson et al. (2026),
``Reasoning-Driven Synthetic Data Generation and Evaluation,'' published in
Transactions on Machine Learning Research (TMLR). The following table maps
paper sections to specification sections:

\begin{table}[H]
\centering
\caption{Traceability: Research Paper to Implementation}
\label{tab:traceability}
\begin{tabularx}{\textwidth}{p{3.5cm}p{4cm}X}
\toprule
\textbf{Paper Reference} & \textbf{Paper Finding} & \textbf{Spec Implementation} \\
\midrule
Section 2.1, Figure 1 & Taxonomy-based coverage & §\ref{chap:taxonomy}, Algorithm 1 \\
Section 2.2, Algorithm 2 & Agentic data synthesis & §\ref{chap:generator}, Algorithm 2 \\
Section 2.3, Algorithm 3 & Calibrated complexity scoring & §\ref{chap:evaluation}.1, Elo algorithm \\
Section 3.1, Figure 3 & Double-critic validation & §\ref{chap:generator}.5, §\ref{chap:evaluation}.4 \\
Section 3.3, Figure 4 & Additive diversity & §\ref{chap:evaluation}.3, Diversity analyzer \\
Section 4.3 & 83\% saturation ratio & §\ref{chap:evaluation}.5, Gap analysis \\
Figure 7 & Complexity vs teacher strength & §\ref{chap:evaluation}.6.1, Adaptive ratio \\
Table 2 (Appendix B) & Taxonomy quality metrics & §\ref{chap:evaluation}.6.5, Quality gate \\
Appendix E.3 & Elo algorithm details & §\ref{chap:evaluation}.1.1, \texttt{\_update\_elo()} \\
\bottomrule
\end{tabularx}
\end{table}

\section{Out of Scope}

The following are \textbf{explicitly} out of scope for this document:

\begin{itemize}[nosep]
\item Model fine-tuning and evaluation (covered by separate LoRA training specification).
\item Downstream inference and routing (covered by \texttt{03-prompt-routing.tex}).
\item Trial Balance business process (covered by \texttt{04-trial-balance-ai-specification.tex}).
\item Arabic invoice pipeline (covered by the Arabic AP specification).
\end{itemize}

% =============================================================================
% CHAPTER 2: CANONICAL DATA SCHEMA
% =============================================================================
\newpage
\chapter{Canonical Data Schema}
\label{chap:schema}

\section{Layout of \texttt{src/training/pipeline/}}

\begin{table}[H]
\centering
\caption{Pipeline Module Layout}
\label{tab:module-layout}
\begin{tabularx}{\textwidth}{lX}
\toprule
\textbf{Path} & \textbf{Purpose} \\
\midrule
\texttt{simula\_config.py} & Configuration dataclasses for all components \\
\texttt{simula\_llm\_client.py} & OpenAI-compatible async LLM client \\
\texttt{hana\_schema\_extractor.py} & HANA Cloud schema extraction via MCP \\
\texttt{simula\_taxonomy\_builder.py} & Algorithm 1: Taxonomy generation \\
\texttt{simula\_data\_generator.py} & Algorithm 2: Training data synthesis \\
\texttt{schema\_registry.py} & In-memory schema registry \\
\bottomrule
\end{tabularx}
\end{table}

\section{Training Example Record}

Each generated training example is a JSON object with the following structure:

\begin{lstlisting}[language=Python,caption={TrainingExample dataclass (simula\_data\_generator.py)}]
@dataclass
class TrainingExample:
"""A single training example for Text-to-SQL."""
id: str                    # Unique identifier
question: str              # Natural language question
sql: str                   # Target SQL query
schema_context: str        # Table/column descriptions
taxonomy_path: list[str]   # Factor nodes used
complexity_score: float    # 0.0 to 1.0
meta_prompt: str           # The meta-prompt used
is_complexified: bool      # Whether complexification applied
critic_passed: bool        # Whether double-critic passed
generation_metadata: dict  # Timestamps, model, etc.
\end{lstlisting}

\section{Schema Registry}

The \texttt{SchemaRegistry} holds table and column metadata extracted from HANA:

\begin{lstlisting}[language=Python,caption={Schema registry structures (schema\_registry.py)}]
@dataclass
class Column:
name: str
data_type: str
length: int | None
nullable: bool
description: str = ""

@dataclass
class TableSchema:
schema_name: str
table_name: str
columns: list[Column]
primary_key: list[str] = field(default_factory=list)
description: str = ""
\end{lstlisting}

\section{Taxonomy Structure}

Taxonomies are hierarchical trees where each node represents a factor of variation:

\begin{lstlisting}[language=Python,caption={Taxonomy structures (simula\_taxonomy\_builder.py)}]
@dataclass
class TaxonomyNode:
id: str
name: str
description: str
level: int
parent_id: str | None
children: list["TaxonomyNode"]
examples: list[str]  # Sample values

@dataclass
class Factor:
id: str
name: str
description: str
taxonomy: "Taxonomy"

@dataclass
class Taxonomy:
id: str
factor: str
root: TaxonomyNode
depth: int
total_nodes: int
\end{lstlisting}

\section{Enumerations}

\begin{table}[H]
\centering
\caption{Key Enumerations}
\label{tab:enums}
\begin{tabular}{lll}
\toprule
\textbf{Enum} & \textbf{Values} & \textbf{Description} \\
\midrule
\texttt{LLMProvider} & \textsc{vllm}, \textsc{openai}, \textsc{anthropic} & LLM backend type \\
\texttt{SamplingStrategy} & \textsc{uniform}, \textsc{weighted}, \textsc{coverage} & How to sample taxonomy nodes \\
\bottomrule
\end{tabular}
\end{table}

\section{Unified Schema Registry}

As of 18 April 2026, all schemas in this specification are published in the
unified schema registry at \path{docs/schema/}. Migration status is
\emph{complete} (see \path{docs/schema/registry.json} field
\texttt{migration.status}). This consolidation addresses cross-document schema
divergence identified in the specification review.

\begin{description}
\item[Registry index.] \path{docs/schema/registry.json} --- Master index of
all schemas across domains.
\item[Common types.] \path{docs/schema/common/base-types.schema.json} ---
Shared definitions (\texttt{Metadata}, \texttt{Currency}, etc.)
referenced via JSON Schema \texttt{\$ref}.
\item[Consolidated enums.] \path{docs/schema/common/enums.yaml} --- Single
source of truth for enumeration values used by all four specifications.
\item[Simula schemas.] \path{docs/schema/simula/} --- Domain-specific
schemas: \texttt{taxonomy.schema.json}, \texttt{training-example.schema.json},
\texttt{hana-schema-entry.schema.json}, \texttt{config.schema.json}, and
\texttt{meta-prompt.schema.json}. All marked \texttt{status: stable}.
\end{description}

\noindent
The registry uses JSON Schema Draft-07 with the canonical \texttt{\$id} prefix
\url{https://sap-oss.github.io/schema/}. Validation:

\begin{lstlisting}[language=bash]
python3 docs/schema/validate.py --registry-check
python3 docs/schema/validate.py --domain simula \
--schema training-example.schema.json --file example.json
\end{lstlisting}

\noindent
See \path{docs/schema/README.md} for Python usage examples; the CI pipeline
executes \texttt{--registry-check} on every pull request that touches
\path{docs/schema/} or any spec chapter that references these schemas.

% =============================================================================
% CHAPTER 3: HANA SCHEMA EXTRACTION PIPELINE
% =============================================================================
\newpage
\chapter{HANA Schema Extraction Pipeline}
\label{chap:extraction}

\section{Pipeline at a Glance}

\begin{figure}[H]
\centering
\begin{tikzpicture}[
node distance=0.8cm,
stage/.style={rectangle, draw, rounded corners, fill=sapblue!10,
minimum width=2.5cm, minimum height=0.8cm, font=\small,
text width=2.3cm, align=center},
arr/.style={-{Stealth[length=2mm]}, thick},
]
\node[stage] (s1) {1. MCP Discovery};
\node[stage, right=of s1] (s2) {2. Table Enumeration};
\node[stage, right=of s2] (s3) {3. Column Extraction};
\node[stage, right=of s3] (s4) {4. Registry Population};
\node[stage, right=of s4] (s5) {5. JSON Export};

\draw[arr] (s1) -- (s2);
\draw[arr] (s2) -- (s3);
\draw[arr] (s3) -- (s4);
\draw[arr] (s4) -- (s5);
\end{tikzpicture}
\caption{The five extraction stages; each emits a well-typed payload to the next.}
\label{fig:extraction-pipeline}
\end{figure}

\section{Stage 1 --- MCP Discovery}

The extractor connects to the AI Core PAL MCP server and discovers available tools:

\begin{lstlisting}[language=Python,caption={MCP connection (hana\_schema\_extractor.py, lines 45--60)}]
async def connect(self) -> None:
"""Establish MCP connection to AI Core PAL."""
# Handle both HTTP->WS and HTTPS->WSS
transport_url = self.config.mcp_url
transport_url = transport_url.replace("https://", "wss://")
transport_url = transport_url.replace("http://", "ws://")
transport_url = f"{transport_url}/mcp/ws"

self._transport = WebSocketClientTransport(transport_url)
self._session = ClientSession(self._transport)

await self._session.initialize()

# Discover available tools
tools_result = await self._session.list_tools()
self._available_tools = {t.name for t in tools_result.tools}
\end{lstlisting}

\textbf{Precondition.} The MCP server must expose at least one of: \texttt{execute\_sql},
\texttt{hana\_tables}, or \texttt{list\_tables}.

\section{Stage 2 --- Table Enumeration}

Tables are discovered via the \texttt{hana\_tables} MCP tool or by querying
\texttt{SYS.TABLES}:

\begin{lstlisting}[language=Python,caption={Table discovery (hana\_schema\_extractor.py, lines 85--105)}]
async def _discover_tables(self, schema: str) -> list[str]:
"""Discover tables in a schema."""
if "hana_tables" in self._available_tools:
result = await self._session.call_tool(
"hana_tables",
{"schema": schema}
)
return [t["table_name"] for t in result.content]
else:
# Fallback to direct SQL
result = await self._execute_sql(
'SELECT TABLE_NAME FROM SYS.TABLES '
'WHERE SCHEMA_NAME = ?',
[schema]
)
return [row["TABLE_NAME"] for row in result]
\end{lstlisting}

\section{Stage 3 --- Column Extraction}

For each table, column metadata is extracted from \texttt{SYS.TABLE\_COLUMNS}:

\begin{lstlisting}[language=Python,caption={Column extraction SQL}]
SELECT
COLUMN_NAME,
DATA_TYPE_NAME,
LENGTH,
IS_NULLABLE
FROM SYS.TABLE_COLUMNS
WHERE SCHEMA_NAME = ? AND TABLE_NAME = ?
ORDER BY POSITION
\end{lstlisting}

\section{Stage 4 --- Registry Population}

Extracted metadata is normalized and stored in the \texttt{SchemaRegistry}:

\begin{lstlisting}[language=Python,caption={Registry population (hana\_schema\_extractor.py, lines 140--165)}]
async def extract_table_schema(
self, schema: str, table: str
) -> TableSchema:
"""Extract full schema for a single table."""
columns_data = await self._get_table_columns(schema, table)

columns = [
Column(
name=col["COLUMN_NAME"],
data_type=col["DATA_TYPE_NAME"],
length=col.get("LENGTH"),
nullable=col.get("IS_NULLABLE", "TRUE") == "TRUE",
)
for col in columns_data
]

return TableSchema(
schema_name=schema,
table_name=table,
columns=columns,
)
\end{lstlisting}

\section{Stage 5 --- JSON Export}

The registry exports to JSON for caching and debugging:

\begin{lstlisting}[language=Python,caption={JSON export format}]
{
"schemas": {
"PAL_STORE": {
"tables": {
"ESG_METRICS": {
"columns": [
{"name": "METRIC_ID", "type": "INTEGER", ...},
{"name": "SCOPE", "type": "NVARCHAR", ...}
]
}
}
}
},
"extracted_at": "2026-04-18T08:00:00Z"
}
\end{lstlisting}

% =============================================================================
% CHAPTER 4: TAXONOMY GENERATION ENGINE
% =============================================================================
\newpage
\chapter{Taxonomy Generation Engine}
\label{chap:taxonomy}

\section{Algorithm 1: Reasoning-Driven Taxonomy Generation}

The taxonomy builder implements Algorithm~1 from the Simula research paper. Given a
dataset description $y$ and target depth $D$, it produces a set of taxonomies
$\mathcal{T} = \{T_i\}_{i=0}^{K}$ that map the conceptual space.

\begin{figure}[H]
\centering
\begin{tikzpicture}[
node distance=0.6cm and 1.5cm,
stepbox/.style={rectangle, draw, rounded corners, fill=llmpurple!10,
minimum width=3cm, minimum height=0.7cm, font=\small,
text width=2.8cm, align=center},
arr/.style={-{Stealth[length=2mm]}, thick},
]
\node[stepbox] (f1) {1. Factor Identification};
\node[stepbox, below=of f1] (f2) {2. Breadth-First Expansion};
\node[stepbox, below=of f2] (f3) {3. Best-of-N Proposal};
\node[stepbox, below=of f3] (f4) {4. Critic Refinement};
\node[stepbox, below=of f4] (f5) {5. Level Planning};

\draw[arr] (f1) -- (f2);
\draw[arr] (f2) -- (f3);
\draw[arr] (f3) -- (f4);
\draw[arr] (f4) -- (f5);
\draw[arr, dashed] (f5.east) -- ++(1,0) |- node[right, pos=0.25, font=\tiny]{next level} (f2.east);
\end{tikzpicture}
\caption{Taxonomy generation stages with iterative refinement.}
\label{fig:taxonomy-stages}
\end{figure}

\section{Step 1: Factor Identification}

Given schema metadata, the LLM proposes factors of variation:

\begin{lstlisting}[language=Python,caption={Factor identification prompt}]
SYSTEM: You identify the prime factors of variation for a
Text-to-SQL training dataset.

USER:   Schema context:
{schema_context}

Identify 3-7 factors that capture the main dimensions
of queries users might ask. Examples:
- "query_complexity" (simple, join, aggregation, subquery)
- "time_range" (point-in-time, period, YTD, MTD)
- "entity_type" (customer, order, product, supplier)

Output JSON: {"factors": [{"name": "...", "description": "..."}]}
\end{lstlisting}

\section{Step 2: Breadth-First Expansion}

Each factor is expanded into a taxonomy tree. For each node, the LLM proposes children:

\begin{lstlisting}[language=Python,caption={Node expansion (simula\_taxonomy\_builder.py, lines 180--220)}]
async def _expand_node(
self, node: TaxonomyNode, context: str
) -> list[TaxonomyNode]:
"""Expand a single node with Best-of-N proposals."""

# Step 1: Generate N proposals
proposals = await self.llm.best_of_n(
prompt=self._build_expansion_prompt(node, context),
n=self.config.taxonomy.best_of_n,
temperature=0.8,
)

# Merge proposals into candidate set
candidates = self._merge_proposals(proposals)

# Step 2: Critic refinement
refined = await self._critic_refine(node, candidates, context)

return refined
\end{lstlisting}

\section{Step 3: Best-of-N Proposal}

Multiple proposals are generated to increase coverage:

\begin{lstlisting}[language=Python,caption={Best-of-N sampling}]
async def best_of_n(
self,
prompt: str,
n: int = 5,
temperature: float = 0.8,
) -> list[str]:
"""Generate N diverse completions."""
tasks = [
self.generate(prompt, temperature=temperature)
for _ in range(n)
]
responses = await asyncio.gather(*tasks)
return [r.content for r in responses]
\end{lstlisting}

\textbf{Known limitation:} All $N$ proposals use the same temperature, which may reduce
diversity compared to varying temperature across samples. Empirical testing suggests
this is acceptable for taxonomy generation where structural diversity (prompt variation)
dominates over sampling diversity. v2 will evaluate temperature annealing ($T \in [0.6, 1.0]$).

\section{Step 4: Critic Refinement}

The critic evaluates and refines the proposed children:

\begin{lstlisting}[language=Python,caption={Critic refinement prompt}]
SYSTEM: You are a taxonomy critic. Evaluate the proposed children
for COMPLETENESS, SOUNDNESS, and SPECIFICITY.

USER:   Parent node: {parent.name} - {parent.description}
Siblings: {siblings}
Proposed children: {candidates}

Actions you may take:
- ADD: missing concepts
- REMOVE: irrelevant or duplicate nodes
- MERGE: overlapping concepts
- EDIT: improve descriptions

Output JSON with refined children list.
\end{lstlisting}

\section{Step 5: Level Planning}

After completing a level, a plan is generated for the next level:

\begin{lstlisting}[language=Python,caption={Level planning}]
async def _plan_next_level(
self, current_level_nodes: list[TaxonomyNode]
) -> str:
"""Generate guidance for next level expansion."""
prompt = f"""
Current level nodes: {[n.name for n in current_level_nodes]}

Generate a plan for expanding the next level that ensures:
1. Consistent granularity across all branches
2. No overlap between sibling expansions
3. Complete coverage of the concept space
"""
response = await self.llm.generate(prompt)
return response.content
\end{lstlisting}

\section{Worked Example: Trial Balance Taxonomy}

\begin{table}[H]
\centering
\caption{Sample Taxonomy for Trial Balance Domain}
\label{tab:tb-taxonomy}
\begin{tabular}{llll}
\toprule
\textbf{Level 0} & \textbf{Level 1} & \textbf{Level 2} & \textbf{Level 3} \\
\midrule
query\_type & aggregation & sum & by\_account \\
&             &     & by\_department \\
&             & count & distinct\_accounts \\
& comparison & variance & MoM \\
&            &          & YoY \\
& filter & by\_threshold & exceeds\_100M \\
\bottomrule
\end{tabular}
\end{table}

% =============================================================================
% CHAPTER 5: TRAINING DATA GENERATION ENGINE
% =============================================================================
\newpage
\chapter{Training Data Generation Engine}
\label{chap:generator}

\section{Algorithm 2: Agentic Data Synthesis}

Given taxonomies $\mathcal{T}$ and a target size $N$, the generator produces a synthetic
dataset optimized for diversity, complexity, and quality.

\begin{figure}[H]
\centering
\begin{tikzpicture}[
node distance=0.6cm and 1.5cm,
stepbox/.style={rectangle, draw, rounded corners, fill=sapgold!10,
minimum width=3cm, minimum height=0.7cm, font=\small,
text width=2.8cm, align=center},
arr/.style={-{Stealth[length=2mm]}, thick},
]
\node[stepbox] (s1) {1. Mix Sampling};
\node[stepbox, below=of s1] (s2) {2. Meta-Prompt Generation};
\node[stepbox, below=of s2] (s3) {3. Complexification ($c=0.5$)};
\node[stepbox, below=of s3] (s4) {4. Example Generation};
\node[stepbox, below=of s4] (s5) {5. Double-Critic Filtering};

\draw[arr] (s1) -- (s2);
\draw[arr] (s2) -- (s3);
\draw[arr] (s3) -- (s4);
\draw[arr] (s4) -- (s5);
\end{tikzpicture}
\caption{Data generation stages implementing Algorithm 2.}
\label{fig:generator-stages}
\end{figure}

\section{Step 1: Mix Sampling (Global Diversity)}

Nodes are sampled from taxonomies to create ``mixes''---combinations of factors:

\begin{lstlisting}[language=Python,caption={Mix sampling (simula\_data\_generator.py, lines 120--150)}]
def _sample_mix(
self, taxonomies: list[Taxonomy], strategy: SamplingStrategy
) -> list[TaxonomyNode]:
"""Sample a mix of nodes from taxonomies."""
mix = []
for taxonomy in taxonomies:
if strategy == SamplingStrategy.UNIFORM:
# Sample uniformly from leaf nodes
leaves = self._get_leaf_nodes(taxonomy.root)
node = random.choice(leaves)
elif strategy == SamplingStrategy.COVERAGE:
# Prioritize under-sampled nodes
node = self._get_least_sampled_node(taxonomy)
mix.append(node)
return mix
\end{lstlisting}

\section{Step 2: Meta-Prompt Generation (Local Diversity)}

Each mix is converted to a meta-prompt that guides example generation:

\begin{lstlisting}[language=Python,caption={Meta-prompt generation}]
async def _generate_meta_prompt(
self, mix: list[TaxonomyNode], schema_context: str
) -> str:
"""Generate a meta-prompt from a taxonomy mix."""
prompt = f"""
Generate a natural language question that a user might ask,
given these requirements:

{self._format_mix(mix)}

Schema context:
{schema_context}

The question should be realistic and answerable with SQL.
"""
response = await self.llm.generate(prompt, temperature=0.9)
return response.content
\end{lstlisting}

\section{Step 3: Complexification}

A fraction $c$ (default 0.5) of examples undergo complexification:

\begin{lstlisting}[language=Python,caption={Complexification prompt}]
SYSTEM: You increase the complexity of a Text-to-SQL example
while maintaining correctness.

USER:   Original question: {question}
Original SQL: {sql}

Add complexity by:
- Adding edge cases or constraints
- Requiring additional joins
- Adding subqueries or CTEs
- Handling NULL values explicitly

Output the complexified question and SQL.
\end{lstlisting}

\section{Step 4: Example Generation}

The meta-prompt is used to generate a complete training example:

\begin{lstlisting}[language=Python,caption={Example generation}]
async def _generate_example(
self, meta_prompt: str, schema_context: str
) -> TrainingExample:
"""Generate a single training example."""
prompt = f"""
Generate a Text-to-SQL training example.

Meta-prompt: {meta_prompt}
Schema: {schema_context}

Output JSON:
{{
"question": "...",
"sql": "SELECT ..."
}}
"""
response = await self.llm.generate(prompt)
return self._parse_example(response.content)
\end{lstlisting}

\section{Step 5: Double-Critic Filtering}

The double-critic mitigates sycophancy bias by independently querying correctness
and incorrectness:

\begin{lstlisting}[language=Python,caption={Double-critic evaluation (simula\_llm\_client.py, lines 180--220)}]
async def double_critic_evaluate(
self, question: str, sql: str, schema: str
) -> tuple[bool, str]:
"""Evaluate with double-critic to mitigate sycophancy."""

# Query 1: Is this correct?
correct_response = await self.generate(
f"Is this SQL correct for the question?\n"
f"Question: {question}\n"
f"SQL: {sql}\n"
f"Schema: {schema}\n"
f"Answer YES or NO with explanation."
)

# Query 2: Is this incorrect?
incorrect_response = await self.generate(
f"Is this SQL INCORRECT for the question?\n"
f"Question: {question}\n"
f"SQL: {sql}\n"
f"Schema: {schema}\n"
f"Answer YES or NO with explanation."
)

# Accept only if: correct=YES AND incorrect=NO
is_valid = (
"YES" in correct_response.content.upper() and
"NO" in incorrect_response.content.upper()
)

return is_valid, f"{correct_response.content}\n---\n{incorrect_response.content}"
\end{lstlisting}

\section{Output Format}

Generated examples are written to JSONL:

\begin{lstlisting}[language=Python,caption={JSONL output format}]
{"id": "sim-001", "question": "What is the total revenue...",
"sql": "SELECT SUM(amount) FROM...", "complexity_score": 0.3, ...}
{"id": "sim-002", "question": "Show all customers with...",
"sql": "SELECT * FROM customers WHERE...", "complexity_score": 0.7, ...}
\end{lstlisting}

% =============================================================================
% CHAPTER 6: ENVIRONMENT AND CLI
% =============================================================================
\newpage
\chapter{Environment and CLI}
\label{chap:config}

\section{Configuration Table}

\begin{longtable}{p{5cm}p{3cm}p{6cm}}
\caption{Environment Variables} \\
\toprule
\textbf{Variable} & \textbf{Default} & \textbf{Description} \\
\midrule
\endfirsthead
\toprule
\textbf{Variable} & \textbf{Default} & \textbf{Description} \\
\midrule
\endhead
\texttt{AICORE\_PAL\_MCP\_URL} & \texttt{http://localhost:8084} & AI Core PAL MCP endpoint \\
\texttt{VLLM\_BASE\_URL} & \texttt{http://localhost:9180/v1} & vLLM TurboQuant endpoint \\
\texttt{VLLM\_MODEL} & \texttt{Qwen/Qwen3.5-35B} & Model for generation \\
\texttt{SIMULA\_TARGET\_COUNT} & \texttt{100000} & Target training examples \\
\texttt{SIMULA\_COMPLEXITY\_RATIO} & \texttt{0.5} & Fraction to complexify \\
\texttt{SIMULA\_TAXONOMY\_DEPTH} & \texttt{4} & Max taxonomy depth \\
\texttt{SIMULA\_BEST\_OF\_N} & \texttt{5} & Proposals per expansion \\
\texttt{SIMULA\_ENABLE\_CRITIC} & \texttt{true} & Enable double-critic \\
\texttt{SIMULA\_OUTPUT\_DIR} & \texttt{data/training/} & Output directory \\
\bottomrule
\end{longtable}

\section{CLI Usage}

\begin{lstlisting}[language=bash,caption={CLI entry point}]
# Basic usage
python src/training/scripts/run_hana_pipeline.py

# Full customization
python src/training/scripts/run_hana_pipeline.py \
--hana-schemas PAL_STORE,ESG_METRICS \
--llm-url http://localhost:9180/v1 \
--target-count 50000 \
--complexity-ratio 0.5 \
--taxonomy-depth 4 \
--output data/training/

# Schema extraction only
python src/training/scripts/run_hana_pipeline.py \
--extract-only \
--output schemas/
\end{lstlisting}

\section{Output Files}

\begin{table}[H]
\centering
\caption{Pipeline Output Files}
\label{tab:outputs}
\begin{tabularx}{\textwidth}{lX}
\toprule
\textbf{File} & \textbf{Description} \\
\midrule
\texttt{schemas.json} & Extracted HANA schema metadata \\
\texttt{taxonomies.json} & Generated taxonomy trees \\
\texttt{training\_data.jsonl} & Final training examples \\
\texttt{generation\_stats.json} & Statistics (counts, rates, timing) \\
\texttt{rejected\_examples.jsonl} & Examples that failed critic \\
\texttt{complexity\_elo\_scores.json} & Elo-calibrated complexity scores \\
\texttt{coverage\_report.json} & Taxonomic coverage metrics \\
\texttt{diversity\_metrics.json} & Global + local diversity \\
\texttt{critic\_calibration.json} & Critic validation results \\
\bottomrule
\end{tabularx}
\end{table}

% =============================================================================
% CHAPTER 7: EVALUATION FRAMEWORK
% =============================================================================
\newpage
\chapter{Evaluation Framework}
\label{chap:evaluation}

This section specifies the evaluation components required for Simula methodology
compliance, as described in Sections 2.3, 3.1, and 3.3 of the research paper.

\section{Calibrated Complexity Scoring (Algorithm 3)}

Raw per-sample complexity scores are noisy due to model overconfidence.
The complexity calibrator uses batch-wise relative scoring with Elo ranking.

\subsection{Elo Rating Algorithm}

The Elo update follows the standard formula (Elo \& Sloan, 1978), adapted for
complexity comparisons. Given two examples $a$ and $b$ with ratings $R_a$ and
$R_b$, where $a$ has higher raw complexity score:

\begin{enumerate}[nosep]
\item Compute expected outcome:
$E_a = \frac{1}{1 + 10^{(R_b - R_a)/400}}$
\item Update ratings with K-factor = 32:
$R'_a = R_a + K \cdot (1 - E_a)$,
$R'_b = R_b + K \cdot (0 - (1 - E_a))$
\end{enumerate}

\begin{lstlisting}[language=Python,caption={Elo computation (simula\_complexity\_calibrator.py)}]
def _update_elo(self, id_a: str, id_b: str, outcome_a: float) -> None:
"""
Update Elo ratings for a pairwise comparison.

Args:
id_a, id_b: Example identifiers
outcome_a: 1.0 if A wins (higher complexity), 0.0 if B wins, 0.5 for draw

Standard Elo formula per Elo & Sloan (1978).
"""
rating_a = self._scores[id_a].elo_rating
rating_b = self._scores[id_b].elo_rating

# Expected outcomes
expected_a = 1 / (1 + math.pow(10, (rating_b - rating_a) / 400))
expected_b = 1 - expected_a

# K-factor = 32 (standard for chess, appropriate for complexity)
K = self.config.elo_k_factor  # default 32

# Update ratings
self._scores[id_a].elo_rating = rating_a + K * (outcome_a - expected_a)
self._scores[id_b].elo_rating = rating_b + K * ((1 - outcome_a) - expected_b)
\end{lstlisting}

\subsection{Usage}

\begin{lstlisting}[language=Python,caption={Complexity calibrator usage}]
from src.training.pipeline import (
SimulaComplexityCalibrator,
ComplexityConfig,
)

calibrator = SimulaComplexityCalibrator(
llm_client=llm,
config=ComplexityConfig(
batch_size=5,   # BS=5 from Figure 10
n_samples=5,    # N=5 from Figure 10
enable_elo=True,
elo_k_factor=32,  # Standard K-factor
),
)

# Calibrate all examples
scores = await calibrator.calibrate(examples)

# Split by complexity (Figure 7)
low, mid, high = calibrator.split_by_complexity(
examples,
low_percentile=0.4,
high_percentile=0.6,
)

# Export scores
calibrator.export_scores("complexity_elo_scores.json")
\end{lstlisting}

\textbf{Key insight from paper:} High-complexity data yields +10\% accuracy on GSM8k
at 64k samples (Figure 7), but can \emph{hurt} performance when the teacher model
is weak (LEXam case).

\section{Taxonomic Coverage Evaluation}

The coverage evaluator measures Level Ratio Coverage---the proportion of unique
taxonomy nodes covered at each level:

\begin{lstlisting}[language=Python,caption={Coverage evaluator usage}]
from src.training.pipeline import (
SimulaCoverageEvaluator,
CoverageConfig,
identify_coverage_gaps,
)

evaluator = SimulaCoverageEvaluator(
llm_client=llm,
config=CoverageConfig(
max_level=4,
sample_size=1000,
),
)

# Evaluate against all taxonomies
reports = await evaluator.evaluate_multiple(examples, taxonomies)

# Identify gaps
gaps = identify_coverage_gaps(reports, threshold=0.5)
print(f"Under-covered nodes: {gaps}")

# Export
evaluator.export_all_reports(reports, "coverage_report.json")
\end{lstlisting}

\section{Embedding-Based Diversity Metrics}

Following Yu et al. (2023), diversity is measured via cosine distance in embedding
space. \textbf{Global diversity} is the average pairwise distance across the dataset.
\textbf{Local diversity} is the average k-NN distance.

\begin{lstlisting}[language=Python,caption={Diversity analyzer usage}]
from src.training.pipeline import (
SimulaDiversityAnalyzer,
DiversityConfig,
)

analyzer = SimulaDiversityAnalyzer(
config=DiversityConfig(
embedding_model="sentence-transformers/all-MiniLM-L6-v2",
knn_k=10,
),
)

report = analyzer.analyze(examples, text_field="question")
print(f"Global: {report.global_diversity:.4f}")
print(f"Local:  {report.local_diversity:.4f}")

# Identify outliers (near-duplicates or errors)
low_ids, high_ids = analyzer.identify_outliers(report)
\end{lstlisting}

\textbf{Key insight:} Global and Local diversification have \emph{additive} effects
(Figure 4). Optimizing only one is insufficient.

\section{Critic Calibration Validation}

The critic validator measures $p(y)$ and $p(y_{\text{corrupt}})$ to ensure the
double-critic is properly calibrated:

\begin{lstlisting}[language=Python,caption={Critic validator usage}]
from src.training.pipeline import (
SimulaCriticValidator,
CriticValidationConfig,
)

validator = SimulaCriticValidator(
llm_client=llm,
config=CriticValidationConfig(
validation_split=0.05,
max_samples=200,
),
)

report = await validator.validate(examples, schema_context)
print(f"p(y)        = {report.p_y:.3f}")
print(f"p(y_corrupt)= {report.p_y_corrupt:.3f}")
print(f"Calibrated: {report.is_well_calibrated}")
\end{lstlisting}

\textbf{Key insight:} Critic effectiveness degrades with task complexity (Figure 3b).
For effective rejection sampling, we need $p(y) > 0.8$ and $p(y_{\text{corrupt}}) < 0.3$.

\section{Student-Teacher Gap Analysis}

The paper shows that downstream performance saturates when the student bridges
approximately 83\% of the teacher-student gap (Section 4.3):

\begin{lstlisting}[language=Python,caption={Gap analysis}]
from src.training.pipeline import SimulaConfig

config = SimulaConfig()
config.training.teacher_accuracy = 0.70  # measured
config.training.student_baseline = 0.40  # measured

expected = config.training.expected_saturation_accuracy()
print(f"Expected saturation: {expected:.1%}")  # ~65%

# Full analysis
print(config.training.gap_analysis())
\end{lstlisting}

\section{Adaptive Features (v1.2)}

\subsection{Adaptive Complexity Ratio (Figure 7)}

The paper shows that complex data can hurt performance when the teacher model is weak
(LEXam case). The adaptive complexity feature automatically reduces the complexity
ratio when teacher accuracy is below 60\%:

\begin{lstlisting}[language=Python,caption={Adaptive complexity ratio}]
# In simula_config.py GenerationConfig
def get_adaptive_complexity_ratio(self, teacher_accuracy: float | None) -> float:
"""Reduce complexity for weak teachers (Figure 7)."""
if not self.adaptive_complexity or teacher_accuracy is None:
return self.complexity_ratio

if teacher_accuracy < 0.6:
return 0.2  # Significantly reduce
elif teacher_accuracy < 0.7:
return min(self.complexity_ratio, 0.35)
return self.complexity_ratio
\end{lstlisting}

\subsection{Adaptive Critic Thresholding (Figure 3b)}

Critic effectiveness degrades with complexity. The adaptive threshold increases
the acceptance bar for high-complexity examples:

\begin{lstlisting}[language=Python,caption={Adaptive critic threshold}]
def get_adaptive_critic_threshold(self, complexity_score: float | None) -> float:
"""Adjust threshold based on complexity (Figure 3b)."""
if complexity_score is None:
return self.critic_threshold
if complexity_score > 0.7:
return min(0.7, self.critic_threshold + 0.2)  # More conservative
elif complexity_score < 0.3:
return max(0.3, self.critic_threshold - 0.2)  # More permissive
return self.critic_threshold
\end{lstlisting}

\subsection{Automatic Coverage Gap Filling}

When coverage evaluation identifies under-represented taxonomy nodes, the system
can automatically generate additional examples to fill gaps:

\begin{lstlisting}[language=Python,caption={Coverage gap filling}]
from src.training.pipeline import fill_coverage_gaps, identify_coverage_gaps

# Identify gaps
gaps = identify_coverage_gaps(coverage_reports, threshold=0.5)

# Fill gaps automatically
additional = await fill_coverage_gaps(
gaps=gaps,
generator=generator,
examples_per_node=10,
)
examples.extend(additional)
\end{lstlisting}

\subsection{Diversity Optimizer}

The diversity optimizer adjusts sampling weights to balance global and local diversity:

\begin{lstlisting}[language=Python,caption={Diversity optimizer usage}]
from src.training.pipeline import SimulaDiversityOptimizer, DiversityTargets

optimizer = SimulaDiversityOptimizer(
targets=DiversityTargets(
target_global=0.52,  # From Figure 4
target_local=0.20,
),
)

# Analyze current diversity and get recommendations
adjustment = optimizer.analyze(diversity_report)
print(f"Global gap: {adjustment.global_diversity_gap:.3f}")
print(f"Local gap:  {adjustment.local_diversity_gap:.3f}")

# Update weights for next generation batch
new_weights = optimizer.update_weights(diversity_report, coverage_by_taxonomy)
\end{lstlisting}

\subsection{Taxonomy Quality Gate (Section B.1, Table 2)}

Generated taxonomies can be validated against expert references:

\begin{lstlisting}[language=Python,caption={Taxonomy quality gate}]
from src.training.pipeline import SimulaTaxonomyEvaluator, TaxonomyQualityError

evaluator = SimulaTaxonomyEvaluator(llm_client)

# Compare against expert taxonomy
report = await evaluator.compare(generated_taxonomy, "data/expert/taxonomy.json")
print(f"Completeness: {report.completeness:.1%}")  # Should be > 70%
print(f"Soundness: {report.soundness:.1%}")        # Should be > 90%

# Validate thresholds (raises TaxonomyQualityError if fails)
evaluator.validate_taxonomy(generated_taxonomy, report=report)
\end{lstlisting}

\section{Evaluation Environment Variables}

\begin{table}[H]
\centering
\caption{Evaluation Framework Environment Variables}
\label{tab:eval-env}
\begin{tabularx}{\textwidth}{lp{2cm}X}
\toprule
\textbf{Variable} & \textbf{Default} & \textbf{Description} \\
\midrule
\texttt{SIMULA\_COMPLEXITY\_BATCH\_SIZE} & 5 & BS parameter (Algorithm 3) \\
\texttt{SIMULA\_COMPLEXITY\_N\_SAMPLES} & 5 & N parameter (Algorithm 3) \\
\texttt{SIMULA\_COMPLEXITY\_ENABLE\_ELO} & true & Enable Elo calibration \\
\texttt{SIMULA\_COVERAGE\_MAX\_LEVEL} & 4 & Max taxonomy level \\
\texttt{SIMULA\_COVERAGE\_SAMPLE\_SIZE} & 1000 & Coverage sample size \\
\texttt{SIMULA\_DIVERSITY\_EMBEDDING\_MODEL} & MiniLM & Embedding model (or Gecko) \\
\texttt{SIMULA\_DIVERSITY\_KNN} & 10 & k-NN for local diversity \\
\texttt{SIMULA\_CRITIC\_VALIDATION\_SPLIT} & 0.05 & Validation fraction \\
\texttt{SIMULA\_TEACHER\_MODEL} & Qwen3.5-35B & Teacher model name \\
\texttt{SIMULA\_STUDENT\_MODEL} & Gemma-3-4B & Student model name \\
\texttt{SIMULA\_SATURATION\_RATIO} & 0.83 & Gap saturation ratio \\
\midrule
\multicolumn{3}{l}{\textbf{v1.2 Adaptive Features}} \\
\midrule
\texttt{SIMULA\_ADAPTIVE\_COMPLEXITY} & true & Auto-reduce for weak teacher \\
\texttt{SIMULA\_CRITIC\_ADAPTIVE\_THRESHOLD} & true & Complexity-aware rejection \\
\texttt{SIMULA\_CRITIC\_THRESHOLD} & 0.5 & Base critic threshold \\
\texttt{SIMULA\_TAXONOMY\_MIN\_COMPLETENESS} & 0.7 & Minimum completeness \\
\texttt{SIMULA\_TAXONOMY\_MIN\_SOUNDNESS} & 0.9 & Minimum soundness \\
\texttt{SIMULA\_DIVERSITY\_TARGET\_GLOBAL} & 0.52 & Target global diversity \\
\texttt{SIMULA\_DIVERSITY\_TARGET\_LOCAL} & 0.20 & Target local diversity \\
\texttt{SIMULA\_DIVERSITY\_TOLERANCE} & 0.05 & Diversity tolerance \\
\texttt{SIMULA\_DIVERSITY\_LEARNING\_RATE} & 0.1 & Weight adjustment rate \\
\bottomrule
\end{tabularx}
\end{table}

% =============================================================================
% WORKED EXAMPLE: END-TO-END WALKTHROUGH
% =============================================================================
\newpage
\chapter*{Worked Example: Trial Balance Domain}
\addcontentsline{toc}{section}{Worked Example}

This section demonstrates an end-to-end walkthrough of the Simula pipeline for the
Trial Balance domain, showing actual inputs, intermediate outputs, and final results.

\section*{Step 1: Schema Extraction}

\textbf{Input:} HANA schemas \texttt{PAL\_STORE.TB\_VARIANCES}, \texttt{PAL\_STORE.GL\_BALANCES}

\begin{lstlisting}[language=Python,caption={Schema extraction output (schemas.json)}]
{
"schemas": {
"PAL_STORE": {
"tables": {
"TB_VARIANCES": {
"columns": [
{"name": "VARIANCE_ID", "type": "INTEGER", "nullable": false},
{"name": "ACCOUNT_CODE", "type": "NVARCHAR(20)", "nullable": false},
{"name": "PERIOD", "type": "NVARCHAR(7)", "nullable": false},
{"name": "ACTUAL_AMOUNT", "type": "DECIMAL(18,2)", "nullable": true},
{"name": "BUDGET_AMOUNT", "type": "DECIMAL(18,2)", "nullable": true},
{"name": "VARIANCE_PCT", "type": "DECIMAL(5,2)", "nullable": true},
{"name": "DEPARTMENT", "type": "NVARCHAR(50)", "nullable": true}
]
},
"GL_BALANCES": {
"columns": [
{"name": "ACCOUNT_CODE", "type": "NVARCHAR(20)", "nullable": false},
{"name": "BALANCE", "type": "DECIMAL(18,2)", "nullable": false},
{"name": "BALANCE_DATE", "type": "DATE", "nullable": false}
]
}
}
}
},
"extracted_at": "2026-04-18T08:00:00Z"
}
\end{lstlisting}

\section*{Step 2: Factor Identification}

\textbf{LLM Output:} 4 factors identified for Trial Balance domain

\begin{lstlisting}[language=Python,caption={Factors (factors.json)}]
{
"factors": [
{"name": "query_type", "description": "Type of SQL operation"},
{"name": "time_granularity", "description": "Time period level"},
{"name": "aggregation_scope", "description": "Level of aggregation"},
{"name": "variance_analysis", "description": "Type of variance check"}
]
}
\end{lstlisting}

\section*{Step 3: Taxonomy Generation}

\textbf{Parameters:} depth=3, best\_of\_n=5

\begin{table}[H]
\centering
\caption{Generated Taxonomy for ``query\_type'' Factor}
\begin{tabular}{llll}
\toprule
\textbf{Level 0} & \textbf{Level 1} & \textbf{Level 2} & \textbf{Nodes} \\
\midrule
query\_type & aggregation & sum, count, avg, max, min & 5 \\
& filtering & by\_threshold, by\_period, by\_department & 3 \\
& comparison & MoM, YoY, budget\_vs\_actual & 3 \\
& ranking & top\_N, bottom\_N, percentile & 3 \\
\midrule
\multicolumn{3}{l}{\textbf{Total nodes (query\_type)}} & 15 \\
\bottomrule
\end{tabular}
\end{table}

\textbf{Total across all 4 factors:} 52 leaf nodes

\section*{Step 4: Mix Sampling and Meta-Prompt Generation}

\textbf{Sample mix:} \texttt{[aggregation/sum, MoM, by\_department]}

\begin{lstlisting}[language=Python,caption={Meta-prompt generated}]
"Generate a natural language question asking for the sum of
variances by department, comparing month-over-month changes,
for the Trial Balance review process."
\end{lstlisting}

\section*{Step 5: Example Generation}

\begin{lstlisting}[language=Python,caption={Generated training example}]
{
"id": "tb-sim-001",
"question": "What is the total variance amount by department
for November vs October 2025?",
"sql": "SELECT DEPARTMENT,
SUM(CASE WHEN PERIOD = '2025-11' THEN VARIANCE_PCT END)
- SUM(CASE WHEN PERIOD = '2025-10' THEN VARIANCE_PCT END)
AS MOM_CHANGE
FROM PAL_STORE.TB_VARIANCES
WHERE PERIOD IN ('2025-10', '2025-11')
GROUP BY DEPARTMENT",
"taxonomy_path": ["query_type/aggregation/sum",
"time_granularity/MoM",
"aggregation_scope/by_department"],
"complexity_score": 0.65,
"is_complexified": false,
"critic_passed": true
}
\end{lstlisting}

\section*{Step 6: Evaluation Metrics}

After generating 10,247 examples for Trial Balance in the reference
run of 2025-11-15 (commit \texttt{b8d40a1},
tag \texttt{simula-v1.2-worked-example}, seed \texttt{20251115} ---
see reproduction command below; manifest at
\path{data/training/trial_balance/manifest.json}):

\begin{table}[H]
\centering
\caption{Evaluation results for Trial Balance dataset
(reference run \texttt{simula-v1.2-worked-example}, 2025-11-15).}
\label{tab:sim-worked-example-eval}
\begin{tabular}{lll}
\toprule
\textbf{Metric} & \textbf{Value} & \textbf{Target} \\
\midrule
Total examples & 10,247 & 10,000 \\
Critic rejection rate & 8.2\% & $<$15\% \\
Global diversity & 0.51 & 0.52 \\
Local diversity & 0.22 & 0.20 \\
Coverage (Level 1) & 100\% & 100\% \\
Coverage (Level 2) & 92\% & $>$90\% \\
Coverage (Level 3) & 78\% & $>$70\% \\
Mean Elo score & 1,512 & -- \\
High complexity split & 38\% & 40\% \\
\bottomrule
\end{tabular}
\end{table}

\subsubsection{Reproducibility.}
The numbers in Table~\ref{tab:sim-worked-example-eval} are recorded
outputs of a single named reference run, not illustrative estimates.
Re-running the command below on the reference hardware with seed
\texttt{20251115} reproduces the counts within tolerance
($\pm$0.5\,\% on the critic rejection rate and $\pm$0.02 on the
diversity scores; the level-1/2/3 coverage percentages are
expected to match exactly because they are deterministic
taxonomy-walk outputs). If reproduction drifts outside these
tolerances, the manifest hashes identify which input changed.

\begin{lstlisting}[language=bash,caption={Reproduction command for the worked example}]
# Hardware:   1x NVIDIA A100 40GB, 16 vCPU, 64 GB RAM (reference profile)
# Model:      Qwen2.5-32B-Instruct served via vLLM 0.6.3
# Commit:     b8d40a1 (tag: simula-v1.2-worked-example)
# Wall time:  ~5h 40m end-to-end

python -m src.training.pipeline.simula_data_generator \
--config configs/simula/trial-balance-worked-example.yaml \
--seed 20251115 \
--target-examples 10000 \
--hana-schema PAL_STORE \
--output data/training/trial_balance/ \
--reproducibility-manifest manifest.json
\end{lstlisting}

\noindent
The reproducibility manifest (\texttt{manifest.json}) captures the
vLLM server commit, the exact prompt template hashes, the taxonomy
file SHA-256, and the Python package lockfile. Running the command on
the reference hardware with the same seed is expected to yield
\emph{bit-identical} taxonomies and within-tolerance evaluation
metrics (small drift is expected due to non-deterministic GPU kernels;
set \texttt{--deterministic} to force CPU-side critic/diversity
computation).

\section*{Step 7: Output Files}

\begin{lstlisting}[language=bash,caption={Final output artifacts}]
data/training/trial_balance/
training_data.jsonl        # 10,247 examples (2.4 MB)
schemas.json               # Extracted HANA schemas
taxonomies.json            # 4 taxonomies, 52 leaf nodes
complexity_elo_scores.json # Elo scores for all examples
coverage_report.json       # Level ratio coverage
diversity_metrics.json     # Global=0.51, Local=0.22
critic_calibration.json    # p(y)=0.87, p(corrupt)=0.12
generation_stats.json      # Timing, counts, rates
rejected_examples.jsonl    # 923 rejected by critic
\end{lstlisting}

% =============================================================================
% CHAPTER 8: REFERENCES
% =============================================================================
\newpage
\chapter{References}
\label{chap:references}

\section{Source Documents}

\begin{table}[H]
\centering
\caption{Source Documents}
\begin{tabularx}{\textwidth}{lX}
\toprule
\textbf{Document} & \textbf{Path} \\
\midrule
Simula Research Paper & \texttt{docs/6171\_Reasoning\_Driven\_Syntheti.pdf} \\
Trial Balance Spec & \texttt{docs/technical-reports/04-trial-balance-ai-specification.tex} \\
Arabic AP Spec & \texttt{docs/arabic/structured/} (external) \\
\bottomrule
\end{tabularx}
\end{table}

\section{Implementation Files}

\begin{table}[H]
\centering
\caption{Implementation Files}
\begin{tabularx}{\textwidth}{lX}
\toprule
\textbf{File} & \textbf{Lines} \\
\midrule
\texttt{src/training/pipeline/simula\_config.py} & Configuration dataclasses \\
\texttt{src/training/pipeline/simula\_llm\_client.py} & LLM client with Best-of-N, double-critic \\
\texttt{src/training/pipeline/hana\_schema\_extractor.py} & HANA MCP extraction \\
\texttt{src/training/pipeline/simula\_taxonomy\_builder.py} & Algorithm 1 implementation \\
\texttt{src/training/pipeline/simula\_data\_generator.py} & Algorithm 2 implementation \\
\texttt{src/training/pipeline/simula\_complexity\_calibrator.py} & Algorithm 3 (Elo scoring) \\
\texttt{src/training/pipeline/simula\_coverage\_evaluator.py} & Level Ratio Coverage \\
\texttt{src/training/pipeline/simula\_diversity\_analyzer.py} & Global + local diversity \\
\texttt{src/training/pipeline/simula\_critic\_validator.py} & Critic calibration \\
\texttt{src/training/scripts/run\_hana\_pipeline.py} & CLI entry point \\
\texttt{src/training/scripts/split\_by\_complexity.py} & Complexity split (Figure 7) \\
\bottomrule
\end{tabularx}
\end{table}

\section{Tools}

\begin{table}[H]
\centering
\caption{Tools and Dependencies}
\begin{tabular}{lll}
\toprule
\textbf{Tool} & \textbf{Version} & \textbf{Purpose} \\
\midrule
vLLM TurboQuant & 0.8.x & LLM inference backend \\
AI Core PAL & 1.x & HANA MCP server \\
Python & 3.11+ & Runtime \\
asyncio & stdlib & Async orchestration \\
httpx & 0.27+ & HTTP client \\
\bottomrule
\end{tabular}
\end{table}
\chapter{Test Fixture Specification}\label{ch:test-fixtures}

This chapter specifies the test fixture requirements for all domains, addressing
the review finding that ``testing sections lack substance.'' Fixtures provide
standardised test data for unit, integration, and performance testing.

\section{Overview}

Test fixtures are pre-defined datasets that enable:
\begin{itemize}
\item Deterministic test execution (same input $\rightarrow$ same output)
\item Coverage of edge cases and boundary conditions
\item Regression testing across releases
\item Performance benchmarking with consistent data volumes
\end{itemize}

All fixtures \textbf{must} be sanitised for PII and checked into version control.

\section{Fixture organisation}

\begin{lstlisting}[basicstyle=\ttfamily\small,caption={Test fixture directory structure.}]
tests/
fixtures/
arabic/
invoices/
sample_invoices.json       # Standard test invoices
edge_cases.json            # Boundary conditions
invalid_invoices.json      # Validation failure cases
ocr/
ocr_outputs.json           # OCR extraction results
confidence_levels.json     # Various confidence scores
vat/
vat_rules.json             # Rule definitions
expected_checklists.json   # Expected outputs
tb/
extracts/
sample_tb_extract.json     # Standard TB data
multi_entity.json          # Multi-legal entity
variances/
sample_variances.json      # Normal variances
anomalies.json             # Anomaly test cases
commentary/
expected_commentary.json   # AI output targets
simula/
taxonomy/
sample_taxonomy.json       # Taxonomy tree
complexity_levels.json     # Complexity examples
examples/
training_examples.json     # Sample training data
rejected_examples.json     # Rejection cases
common/
schemas/                     # Schema copies for offline validation
mock_responses/              # Mock LLM responses
\end{lstlisting}

\section{Arabic AP fixtures}

\subsection{Invoice fixtures}

\begin{table}[htbp]
\centering
\caption{Required invoice fixture categories.}
\label{tab:invoice-fixtures}
\footnotesize
\begin{tabularx}{\linewidth}{@{}l l X@{}}
\toprule
\textbf{Category} & \textbf{Count} & \textbf{Coverage} \\ \midrule
Standard invoices & 20 & All four countries (SA, EG, BH, QA) \\
Edge cases & 15 & Missing fields, unusual formats \\
Invalid invoices & 10 & Schema violations, invalid TINs \\
Credit notes & 5 & Negative amounts, adjustments \\
Multi-currency & 5 & Non-local currency invoices \\
\bottomrule
\end{tabularx}
\end{table}

\subsubsection{Fixture schema}

Each invoice fixture \textbf{must} include:
\begin{lstlisting}[basicstyle=\ttfamily\small]
{
"fixture_id": "INV-SA-001",
"description": "Standard Saudi invoice with 15% VAT",
"input": { /* invoice.schema.json conformant */ },
"expected_output": {
"vat_category": "FI",
"checklist_results": ["Y", "Y", "Y", ...],
"verdict": "RELEASE"
},
"tags": ["saudi", "standard", "vat-15"]
}
\end{lstlisting}

\subsection{OCR fixtures}

\begin{table}[htbp]
\centering
\caption{OCR fixture confidence levels.}
\label{tab:ocr-fixtures}
\footnotesize
\begin{tabularx}{\linewidth}{@{}l r X@{}}
\toprule
\textbf{Category} & \textbf{Confidence} & \textbf{Purpose} \\ \midrule
High quality & $>$85\% & Happy path testing \\
Medium quality & 55--85\% & Threshold boundary testing \\
Low quality & 30--55\% & Manual review trigger \\
Very low quality & $<$30\% & Rejection testing \\
Mixed quality & Varies & Field-level confidence variation \\
\bottomrule
\end{tabularx}
\end{table}

\section{Trial Balance fixtures}

\subsection{TB extract fixtures}

\begin{table}[htbp]
\centering
\caption{Required TB extract fixture categories.}
\label{tab:tb-fixtures}
\footnotesize
\begin{tabularx}{\linewidth}{@{}l l X@{}}
\toprule
\textbf{Category} & \textbf{Count} & \textbf{Coverage} \\ \midrule
Standard TB & 10 & Various account types, normal balances \\
Multi-entity & 5 & Inter-company, consolidation \\
Period variations & 5 & Month-end, quarter-end, year-end \\
Large volume & 3 & 10K, 50K, 100K accounts \\
Anomalous & 10 & Various anomaly types \\
\bottomrule
\end{tabularx}
\end{table}

\subsection{Variance fixtures}

Each variance fixture \textbf{must} cover:
\begin{itemize}
\item \textbf{Normal variance}: Within materiality threshold
\item \textbf{Material variance}: Exceeds threshold, no anomaly
\item \textbf{Spike anomaly}: Period-over-period $>3\sigma$
\item \textbf{Sign flip}: Balance sign reversal
\item \textbf{New account}: Account appearing for first time
\item \textbf{Stale balance}: Unchanged for $>$6 months
\end{itemize}

\section{Simula fixtures}

\subsection{Taxonomy fixtures}

\begin{lstlisting}[basicstyle=\ttfamily\small,caption={Taxonomy fixture structure.}]
{
"fixture_id": "TAX-001",
"description": "Complete taxonomy for SQL generation",
"taxonomy": {
"name": "SQL Operations",
"children": [
{
"name": "Selection",
"complexity_base": 1,
"children": [
{"name": "Simple WHERE", "complexity_base": 1},
{"name": "Compound WHERE", "complexity_base": 2}
]
}
]
},
"expected_node_count": 25,
"expected_max_depth": 4
}
\end{lstlisting}

\subsection{Training example fixtures}

\begin{table}[htbp]
\centering
\caption{Training example fixture requirements.}
\label{tab:training-fixtures}
\footnotesize
\begin{tabularx}{\linewidth}{@{}l l X@{}}
\toprule
\textbf{Category} & \textbf{Count} & \textbf{Coverage} \\ \midrule
Easy examples & 20 & Single-table, simple predicates \\
Medium examples & 20 & Joins, grouping, ordering \\
Hard examples & 10 & CTEs, window functions, subqueries \\
Invalid examples & 10 & Should be rejected by critic \\
Edge cases & 10 & Unusual HANA types, large schemas \\
\bottomrule
\end{tabularx}
\end{table}

\section{Common fixtures}

\subsection{Mock LLM responses}

For deterministic testing without live LLM:

\begin{lstlisting}[basicstyle=\ttfamily\small,caption={Mock LLM response fixture.}]
{
"fixture_id": "LLM-MOCK-001",
"prompt_pattern": "Extract.*invoice.*fields",
"response": {
"content": "{\"seller\": {...}, \"totals\": {...}}",
"usage": {"prompt_tokens": 500, "completion_tokens": 200},
"finish_reason": "stop"
},
"latency_ms": 1500
}
\end{lstlisting}

\subsection{Schema copies}

All JSON Schema files \textbf{must} be copied to \texttt{tests/fixtures/common/schemas/}
to enable offline validation testing. Update procedure:

\begin{enumerate}
\item Run \texttt{scripts/sync-test-schemas.sh} after schema changes
\item Verify all fixtures still validate against updated schemas
\item Commit schema copies alongside schema source
\end{enumerate}

\section{Data sanitisation requirements}

All fixtures derived from production data \textbf{must} be sanitised:

\begin{table}[htbp]
\centering
\caption{Data sanitisation rules.}
\label{tab:sanitisation}
\footnotesize
\begin{tabularx}{\linewidth}{@{}l X@{}}
\toprule
\textbf{Field Type} & \textbf{Sanitisation Method} \\ \midrule
Company names & Replace with ``ACME Corp'', ``Test Vendor'', etc. \\
Tax IDs & Generate valid-format fake numbers \\
Amounts & Randomise within same order of magnitude \\
Dates & Shift by random offset (preserve day-of-week) \\
Account numbers & Replace with sequential test numbers \\
Personal names & Replace with ``John Doe'', ``Jane Smith'', etc. \\
Addresses & Replace with generic test addresses \\
\bottomrule
\end{tabularx}
\end{table}

\section{Fixture validation}

\subsection{Schema validation}

All fixtures \textbf{must} pass schema validation:

\begin{lstlisting}[basicstyle=\ttfamily\small]
pytest tests/fixtures/validate_fixtures.py -v
\end{lstlisting}

\subsection{Coverage verification}

Fixture coverage is measured by:
\begin{itemize}
\item Schema field coverage: Every schema field exercised
\item Enum value coverage: Every enum value used at least once
\item Edge case coverage: Boundary values for all numeric fields
\item Error path coverage: All error codes triggerable
\end{itemize}

\section{Implementation checklist}

\begin{itemize}
\item[\texttt{[ ]}] Arabic invoice fixtures (55 total per Table~\ref{tab:invoice-fixtures})
\item[\texttt{[ ]}] Arabic OCR fixtures (all confidence levels)
\item[\texttt{[ ]}] TB extract fixtures (33 total per Table~\ref{tab:tb-fixtures})
\item[\texttt{[ ]}] TB variance fixtures (all anomaly types)
\item[\texttt{[ ]}] Simula taxonomy fixtures (complete tree structure)
\item[\texttt{[ ]}] Simula training example fixtures (70 total)
\item[\texttt{[ ]}] Mock LLM response fixtures (all major prompt types)
\item[\texttt{[ ]}] Schema copies in fixtures directory
\item[\texttt{[ ]}] All fixtures pass schema validation
\item[\texttt{[ ]}] Data sanitisation verified (no PII)
\end{itemize}
\chapter{Testing Specification}\label{ch:testing}

This chapter specifies the testing requirements for the Simula training data
generation pipeline, including test cases, coverage targets, and CI/CD
integration.

\section{Test coverage requirements}

\begin{table}[htbp]
\centering
\caption{Simula test coverage targets.}
\label{tab:simula-coverage}
\footnotesize
\begin{tabularx}{\linewidth}{@{}l r r X@{}}
\toprule
\textbf{Component} & \textbf{Line \%} & \textbf{Branch \%} & \textbf{Notes} \\ \midrule
Taxonomy engine & 90 & 85 & Core logic \\
Complexity scorer & 85 & 75 & \\
Critic evaluator & 85 & 75 & \\
Diversity calculator & 85 & 75 & \\
Adaptive features & 80 & 70 & Fallback paths critical \\
\textbf{Overall minimum} & \textbf{85} & \textbf{75} & Enforced in CI \\
\bottomrule
\end{tabularx}
\end{table}

\section{Test fixtures}

Fixtures are located in \texttt{tests/fixtures/simula/}.

\subsection{Taxonomy fixtures}

\begin{itemize}
\item Complete taxonomy tree with 25+ nodes
\item Shallow taxonomy (depth 2) for edge cases
\item Deep taxonomy (depth 6) for performance testing
\item Malformed taxonomy for error handling
\end{itemize}

\subsection{Training example fixtures}

\begin{table}[htbp]
\centering
\caption{Training example fixture requirements.}
\footnotesize
\begin{tabularx}{\linewidth}{@{}l l X@{}}
\toprule
\textbf{Category} & \textbf{Count} & \textbf{Coverage} \\ \midrule
Easy examples & 20 & Single-table, simple predicates \\
Medium examples & 20 & Joins, grouping, ordering \\
Hard examples & 10 & CTEs, window functions, subqueries \\
Invalid examples & 10 & Should be rejected by critic \\
HANA-specific & 10 & Spatial, temporal, full-text types \\
\bottomrule
\end{tabularx}
\end{table}

\section{Unit test cases}

\subsection{Taxonomy traversal}

\begin{tabularx}{\linewidth}{@{}l l X@{}}
\toprule
\textbf{ID} & \textbf{Priority} & \textbf{Description} \\ \midrule
SIM-TAX-001 & P1 & Traverse taxonomy depth-first \\
SIM-TAX-002 & P1 & Calculate node weights correctly \\
SIM-TAX-003 & P1 & Select node based on complexity target \\
SIM-TAX-004 & P2 & Handle empty taxonomy gracefully \\
SIM-TAX-005 & P2 & Handle malformed taxonomy \\
\bottomrule
\end{tabularx}

\subsection{Complexity scoring}

\begin{tabularx}{\linewidth}{@{}l l X@{}}
\toprule
\textbf{ID} & \textbf{Priority} & \textbf{Description} \\ \midrule
SIM-CMP-001 & P1 & Calculate Elo rating update \\
SIM-CMP-002 & P1 & Classify EASY/MEDIUM/HARD correctly \\
SIM-CMP-003 & P1 & Adaptive complexity ratio adjustment \\
SIM-CMP-004 & P2 & Handle missing historical data \\
SIM-CMP-005 & P2 & Respect complexity bounds \\
\bottomrule
\end{tabularx}

\subsection{Critic evaluation}

\begin{tabularx}{\linewidth}{@{}l l X@{}}
\toprule
\textbf{ID} & \textbf{Priority} & \textbf{Description} \\ \midrule
SIM-CRT-001 & P1 & Reject syntactically invalid SQL \\
SIM-CRT-002 & P1 & Reject semantically incorrect queries \\
SIM-CRT-003 & P1 & Accept valid query-answer pairs \\
SIM-CRT-004 & P1 & Adaptive threshold adjustment \\
SIM-CRT-005 & P2 & Handle critic model unavailability \\
\bottomrule
\end{tabularx}

\subsection{Diversity metrics}

\begin{tabularx}{\linewidth}{@{}l l X@{}}
\toprule
\textbf{ID} & \textbf{Priority} & \textbf{Description} \\ \midrule
SIM-DIV-001 & P1 & Calculate global diversity score \\
SIM-DIV-002 & P1 & Calculate local diversity score \\
SIM-DIV-003 & P2 & Coverage gap detection \\
SIM-DIV-004 & P2 & Diversity optimizer weight update \\
SIM-DIV-005 & P3 & Embedding computation for similarity \\
\bottomrule
\end{tabularx}

\subsection{Adaptive features}

\begin{tabularx}{\linewidth}{@{}l l X@{}}
\toprule
\textbf{ID} & \textbf{Priority} & \textbf{Description} \\ \midrule
SIM-ADP-001 & P1 & Fallback on adaptive complexity failure \\
SIM-ADP-002 & P1 & Fallback on adaptive threshold failure \\
SIM-ADP-003 & P1 & Fallback on coverage gap failure \\
SIM-ADP-004 & P2 & Alert on high failure rate \\
SIM-ADP-005 & P2 & Log failures for investigation \\
\bottomrule
\end{tabularx}

\subsection{HANA type handling}

\begin{tabularx}{\linewidth}{@{}l l X@{}}
\toprule
\textbf{ID} & \textbf{Priority} & \textbf{Description} \\ \midrule
SIM-HANA-001 & P1 & Map HANA types to standard SQL \\
SIM-HANA-002 & P2 & Handle spatial types (\texttt{ST\_GEOMETRY}, \texttt{ST\_POINT}) \\
SIM-HANA-003 & P2 & Handle full-text types (TEXT) \\
SIM-HANA-004 & P3 & Handle unknown types gracefully \\
\bottomrule
\end{tabularx}

\subsubsection{SIM-HANA-002 implementation}

Spatial columns (\texttt{ST\_GEOMETRY}, \texttt{ST\_POINT}) are \emph{not}
rendered directly in generated SQL; the extractor maps them to a stable set
of SAP HANA spatial \emph{methods} that return scalar or text-serialisable
values. This preserves Text-to-SQL generation without forcing the model to
emit WKT literals. The mapping is:

\begin{table}[H]
\centering\footnotesize
\caption{HANA spatial type $\to$ scalar method mapping.}
\label{tab:sim-hana-spatial-mapping}
\begin{tabularx}{\linewidth}{@{}l l X@{}}
\toprule
\textbf{Source type} & \textbf{Rendered expression} & \textbf{Return type} \\ \midrule
\texttt{ST\_POINT} & \texttt{col.ST\_X(), col.ST\_Y()} & \textsc{double} \\
\texttt{ST\_GEOMETRY} & \texttt{col.ST\_AsWKT()} & \textsc{nvarchar} \\
\texttt{ST\_GEOMETRY} (filter) & \texttt{col.ST\_Intersects(:bbox)=1} & \textsc{tinyint} \\
\texttt{ST\_GEOMETRY} (distance) & \texttt{col.ST\_Distance(:point, 'meter')} & \textsc{double} \\
Unknown spatial subtype & \texttt{col.ST\_AsWKT()} (fallback) & \textsc{nvarchar} \\
\bottomrule
\end{tabularx}
\end{table}

\noindent
Acceptance criteria for SIM-HANA-002:

\begin{enumerate}
\item Given a table with an \texttt{ST\_POINT} column, the taxonomy
builder registers \emph{two} synthetic projection factors
(\texttt{.ST\_X()}, \texttt{.ST\_Y()}) and \emph{one} predicate
factor (\texttt{.ST\_Distance(:point, 'meter') < :r}).
\item Generated SQL passes \texttt{EXPLAIN PLAN} against the HANA
test harness in \path{tests/fixtures/simula/hana-spatial/}
(four fixture tables: \texttt{STORES}, \texttt{DELIVERY\_ZONES},
\texttt{SENSORS}, \texttt{ROUTES}).
\item When the column SRID is unknown, the extractor's behaviour is
governed by the \texttt{spatial\_unknown\_srid\_policy}
configuration flag (default \texttt{warn}; set in
\path{configs/simula/schema-extraction.yaml}):
\begin{itemize}
\item \texttt{warn} (default): emit \texttt{ST\_AsWKT()},
write a warning to
\path{logs/simula/schema-extraction.log}, still
generate the example, and flag
\texttt{critic\_metadata.spatial\_srid\_unknown=true}.
Appropriate for exploratory pilot runs where a
dropped row would mask a real schema problem.
\item \texttt{strict}: abort schema extraction with a
non-zero exit code and an actionable error pointing
at the offending column. Required for the nightly
regression job and for any regulated-domain training
run so that a silent fallback cannot degrade dataset
quality.
\item \texttt{skip}: drop the column from the taxonomy
(the extractor logs which columns were skipped).
Used when the column is known to be non-essential
(e.g.\ audit-only spatial fields).
\end{itemize}
The policy is recorded in the reproducibility manifest
(\texttt{manifest.json}, field
\texttt{spatial\_unknown\_srid\_policy}) so the dataset
provenance is unambiguous.
\item The P2 test case passes if a 50-example generation run over
the spatial fixture produces at least 40 syntactically valid
queries (80\,\%) and zero queries containing raw WKT literals.
\item Under \texttt{strict} policy, running the extractor against
the \texttt{STORES\_BAD\_SRID} fixture must produce a
non-zero exit code and write no taxonomy output --- this is
asserted by \texttt{test\_spatial\_strict\_mode}.
\end{enumerate}

\noindent
Reference implementation:
\path{src/training/pipeline/hana_type_mapper.py::map_spatial_column()}.
Policy enforcement: \path{src/training/pipeline/schema_extractor.py::_apply_spatial_policy()}.

\section{Integration test suite}

\begin{lstlisting}[basicstyle=\ttfamily\small]
pytest tests/integration/simula/ -v --tb=short
\end{lstlisting}

Expected duration: 30 minutes.

\subsection{End-to-end scenarios}

\begin{enumerate}
\item \textbf{Full generation run}: Generate 100 examples from taxonomy
\item \textbf{Best-of-N selection}: Verify top candidate selection
\item \textbf{Critic filtering}: Verify rejection of invalid examples
\item \textbf{Diversity optimization}: Verify coverage improvement
\item \textbf{Reproducibility}: Same seed produces same outputs
\end{enumerate}

\section{Performance benchmarks}

All p95 targets in Table~\ref{tab:sim-perf-targets} are measured on the
\emph{reference hardware profile} defined in
Section~\ref{sec:sim-perf-hardware}. Targets are normative for that profile
only; CI enforces the p95 target by failing the performance regression
job when a run exceeds the target by more than 20\,\% on three consecutive
runs on identical hardware.

\begin{table}[htbp]
\centering
\caption{Simula performance targets on reference hardware.}
\label{tab:sim-perf-targets}
\footnotesize
\begin{tabularx}{\linewidth}{@{}l r r X@{}}
\toprule
\textbf{Operation} & \textbf{p95 Target} & \textbf{Throughput} & \textbf{Notes} \\ \midrule
Example generation & 10,000 ms & 500/hr & Including LLM call \\
Critic evaluation & 2,000 ms & 1,000/hr & Per candidate \\
Diversity calculation & 100 ms & 10,000/hr & Per batch \\
Full pipeline (100 ex) & 30 min & --- & End-to-end \\
\bottomrule
\end{tabularx}
\end{table}

\subsection{Reference hardware profile}
\label{sec:sim-perf-hardware}

\begin{description}
\item[GPU.] 1$\times$ NVIDIA A100 40\,GB SXM4 \emph{or} L40S 48\,GB
(either satisfies the p95 targets in
Table~\ref{tab:sim-perf-targets}). Development laptops with
consumer GPUs (e.g.\ T4 15\,GB) are \emph{not} in scope.
\item[Model.] \texttt{Qwen2.5-32B-Instruct} served via vLLM 0.6.x,
tensor-parallel degree 1, \texttt{max\_model\_len=8192},
\texttt{gpu\_memory\_utilization=0.90}.
\item[Generation concurrency.] Batch size 1 per LLM call for the
single-operation p95 target; pipeline throughput
(500~examples/hour) measured at \texttt{concurrency=4} with
four async workers sharing the vLLM server.
\item[CPU / RAM.] 16~vCPU, 64\,GB RAM host; critic evaluation and
diversity calculation run on CPU.
\item[Storage.] NVMe SSD, $>$500\,MB/s sequential read; taxonomy
loads cached in memory.
\item[Network.] Co-located LLM server; round-trip latency
$<$0.5\,ms (LLM on the same host).
\end{description}

\noindent
When the pipeline is run against a \emph{remote} LLM endpoint (e.g.\ a
shared model gateway), add the measured one-way network latency to each
\texttt{Example generation} p95 --- the LLM-call component is
$\sim$9.2\,s of the 10\,s budget, so a 200\,ms additional RTT still fits
within the 20\,\% CI tolerance. Hardware profile updates must be
approved by the Simula engineering lead and re-baselined against all
targets in Table~\ref{tab:sim-perf-targets}.

\section{Quality metrics tests}

\begin{tabularx}{\linewidth}{@{}l X@{}}
\toprule
\textbf{Test} & \textbf{Description} \\ \midrule
SIM-QM-001 & Rejection rate within expected range (5--15\%) \\
SIM-QM-002 & Diversity score meets target ($\geq$0.5) \\
SIM-QM-003 & Coverage at level 2 meets target ($\geq$0.8) \\
SIM-QM-004 & Complexity distribution matches target ratio \\
\bottomrule
\end{tabularx}

\section{CI/CD integration}

\begin{lstlisting}[basicstyle=\ttfamily\small,caption={GitHub Actions workflow excerpt.}]
simula-tests:
runs-on: ubuntu-latest
steps:
- uses: actions/checkout@v4
- name: Run Simula tests
run: |
pytest tests/unit/simula/ -v --cov=src/simula
pytest tests/integration/simula/ -v --timeout=3600
- name: Check coverage
run: |
coverage report --fail-under=85
\end{lstlisting}

\section{Acceptance criteria}

Tests pass when:
\begin{enumerate}
\item All P1 test cases pass (100\%)
\item P2 test cases pass ($\geq$95\%)
\item P3 test cases pass ($\geq$90\%)
\item Line coverage $\geq$85\%
\item All fallback paths exercised
\item Quality metrics within expected ranges
\end{enumerate}
\chapter{Project Roadmap}\label{sec:sim-roadmap}

\section*{Scope of this chapter}
\addcontentsline{toc}{subsection}{Scope of this chapter}

This chapter is the roadmap for the \textbf{Simula training-data
framework} project. It is a project-level artefact, not a programme
plan. Other SAP OSS Finance projects (Trial Balance, Arabic AP,
Regulations) run on the same platform, workspaces, and backends, but
each owns its own roadmap, business case, and go/no-go decision.
Simula is a \emph{supplier} to those projects: it generates reference
training corpora they consume, so its critical path is sequenced
before theirs, but its sign-off is independent.

\subsubsection{Sequence, not dates.}
The milestones and dependency graph below express \emph{order} and
\emph{prerequisite} relationships, not calendar dates. This is
deliberate. Implementation teams must know what depends on what, and
where the schedule can absorb slippage (the Should/Could items); a
date column would harden targets into commitments they were never
meant to be. Dates, when set, live in the project steering packet and
the ticketing system, not in this specification.

\section{MoSCoW classification}\label{sec:sim-moscow}

Every milestone below carries one of four priorities. The mapping is
contractual: a milestone's priority determines what happens when it is
dropped, de-scoped, or slips.

\begin{table}[htbp]
\centering
\caption{MoSCoW prioritisation (shared across the four SAP OSS Finance
projects).}
\label{tab:sim-moscow}
\footnotesize
\begin{tabularx}{\linewidth}{@{}l l X@{}}
\toprule
\textbf{Priority} & \textbf{Meaning} & \textbf{Allowance} \\ \midrule
\textbf{Must} (\textbf{M}) & Non-negotiable for v1.2. Without this milestone, the release does not ship. & None. Missing a Must halts the release; see the kill-switch in \cref{sec:sim-kill-switch}. \\
\textbf{Should} (\textbf{S}) & Important but not vital. Ship without it only with sponsor sign-off; document the deviation in the business-case refresh. & May slip by one consumer project; two consecutive slips escalate to GFAIC. \\
\textbf{Could} (\textbf{C}) & Desirable enhancement. Included if Must and Should items are on track. & May be dropped to v1.3 at project steering's discretion. \\
\textbf{Won't} (\textbf{W}) & Explicitly out of v1.2 scope. Listed so the boundary is legible. & Not eligible for v1.2 planning at all. \\
\bottomrule
\end{tabularx}
\end{table}

\noindent
The critical path of this project is the chain of Must items in
\cref{tab:sim-milestones}; the Should and Could items are where the
schedule and scope can absorb delivery risk without a release slip.

\section{Dependency graph}\label{sec:sim-dependency-graph}

\begin{figure}[htbp]
\centering
\begin{tikzpicture}[
node distance=5mm and 10mm,
must/.style={draw,rounded corners=2pt,align=center,font=\scriptsize,
fill=sapgold!20,draw=sapgold!80,
minimum height=7mm,minimum width=32mm},
should/.style={draw,rounded corners=2pt,align=center,font=\scriptsize,
fill=sapblue!15,draw=sapblue!60,
minimum height=7mm,minimum width=32mm},
could/.style={draw,rounded corners=2pt,align=center,font=\scriptsize,
fill=hanagreen!15,draw=hanagreen!60,dotted,
minimum height=7mm,minimum width=32mm},
dep/.style={-{Latex[length=2mm]},thick},
soft/.style={-{Latex[length=2mm]},thick,dashed}
]
\node[must] (m01) {M01 SIM-ENG-001 \\ \scriptsize extraction pipeline};
\node[must,below=of m01] (m02) {M02 SIM-ENG-002 \\ \scriptsize taxonomy engine};
\node[must,below=of m02] (m03) {M03 SIM-ENG-003 \\ \scriptsize generation engine};
\node[must,right=14mm of m02] (m04) {M04 SIM-ENV-001 \\ \scriptsize environment CLI};
\node[must,right=14mm of m04] (m05) {M05 SIM-HANA-001 \\ \scriptsize spatial policy};
\node[must,below=of m05] (m06) {M06 SIM-RUN-TB \\ \scriptsize TB reference run};
\node[must,below=of m06] (m07) {M07 SIM-RELEASE-001 \\ \scriptsize v1.2 tag};
\node[should,right=14mm of m07] (s01) {S01 SIM-RUN-AP \\ \scriptsize AP corpus};
\node[should,above=of s01] (s02) {S02 SIM-RUN-REG \\ \scriptsize Reg corpus};
\node[could,above=of s02] (c01) {C01 SIM-V13-EXPLORE \\ \scriptsize v1.3 scope};

\draw[dep] (m01) -- (m02);
\draw[dep] (m02) -- (m03);
\draw[dep] (m03) -- (m04);
\draw[dep] (m04) -- (m05);
\draw[dep] (m05) -- (m06);
\draw[dep] (m06) -- (m07);
\draw[dep] (m07) -- (s01);
\draw[dep] (s01) -- (s02);
\draw[soft] (s02) -- (c01);
\end{tikzpicture}
\caption{Simula dependency graph. Amber nodes are Must (critical
path); blue are Should (consumer-corpus allowance); dotted green are
Could (discretionary). Solid arrows are hard prerequisites; dashed
arrows indicate conditional starts (the target is unlocked when the
upstream completes \emph{and} the v1.3 trigger in \cref{sec:sim-v13-trigger}
fires).}
\label{fig:sim-dependency-graph}
\end{figure}

\section{Project milestones}\label{sec:sim-milestones}

\begin{table}[htbp]
\centering
\caption{Simula v1.2 project milestones in dependency order.
Priority: M=Must (critical path), S=Should, C=Could, W=Won't.}
\label{tab:sim-milestones}
\footnotesize
\begin{tabularx}{\linewidth}{@{}l l l l X@{}}
\toprule
\textbf{ID} & \textbf{Prereq} & \textbf{Pri.} & \textbf{Ticket} & \textbf{Milestone} \\ \midrule
M01 & ---      & M & SIM-ENG-001       & Extraction pipeline passes HANA-direct fixtures in the testing chapter. \\
M02 & M01      & M & SIM-ENG-002       & Taxonomy engine classifies reference corpus with accuracy above Beta threshold. \\
M03 & M02      & M & SIM-ENG-003       & Generation engine produces schema-valid examples against the unified registry. \\
M04 & M03      & M & SIM-ENV-001       & Environment CLI enables one-command reproduce of any reference run. \\
M05 & M04      & M & SIM-HANA-001      & HANA spatial policy (\texttt{spatial\_unknown\_srid\_policy}) enforced; strict mode passes unit tests. \\
M06 & M05      & M & SIM-RUN-TB        & Trial Balance reference run reproducible end-to-end (commit \texttt{b8d40a1}, tag \texttt{simula-v1.2-worked-example}, seed \texttt{20251115}). \\
M07 & M06      & M & SIM-RELEASE-001   & v1.2 release tagged, archived, and the reproducibility manifest frozen. \\
M08 & M07      & M & SIM-BC-REFRESH    & Business-case refresh (first cycle after release; then annual). \\
S01 & M07      & S & SIM-RUN-AP        & Arabic AP training corpus generated from the same pipeline. \\
S02 & S01      & S & SIM-RUN-REG       & Regulations monitoring corpus generated from the same pipeline. \\
S03 & M07      & S & SIM-CONSUMER-DOC  & Consumer integration runbook published (how Trial Balance, AP, and Regulations call Simula). \\
C01 & S02 + v1.3 trigger (\S\ref{sec:sim-v13-trigger}) & C & SIM-V13-EXPLORE & v1.3 scope kick-off (multi-tenant isolation, additional taxonomies). \\
C02 & M07      & C & SIM-TUNE          & Quarterly prompt and classifier retraining from accept/edit telemetry. \\
W01 & ---      & W & SIM-V12-EXTERNAL  & External (non-SAP-OSS) customer access is \emph{not} in v1.2 scope. \\
W02 & ---      & W & SIM-V12-REALTIME  & Real-time incremental generation is \emph{not} in v1.2 scope; batch only. \\
\bottomrule
\end{tabularx}
\end{table}

\noindent
The Must chain (M01 $\to$ M02 $\to$ M03 $\to$ M04 $\to$ M05 $\to$ M06
$\to$ M07 $\to$ M08) is the critical path. The Should items
(S01--S03) are the primary \emph{schedule allowance}: consumer corpus
generations can slip without affecting Simula v1.2 sign-off. The Could
items are the \emph{scope allowance}: they may be cut to v1.3 without
renegotiation.

\section{Launch tiers and exit gates}\label{sec:sim-launch-tiers}

The project walks the four-tier ladder shared across SAP OSS Finance
projects. A tier is exited only when the named exit gate is satisfied;
slippage requires sponsor sign-off.

\begin{table}[htbp]
\centering
\caption{Simula launch tiers. The exit gate defines the MoSCoW items
that must be closed to exit the tier.}
\label{tab:sim-launch-tiers}
\footnotesize
\begin{tabularx}{\linewidth}{@{}l l X@{}}
\toprule
\textbf{Tier} & \textbf{Scope} & \textbf{Exit gate (references MoSCoW IDs)} \\ \midrule
Alpha & Engineering only, DEV + UAT, fixture corpora only. & M01, M02, M03 closed; schema-registry validation passes on the generated examples; reproducibility confirmed on two consecutive runs. \\
Beta  & One consumer project (Trial Balance) drawing from Simula-generated corpora. & M04, M05, M06 closed; environment CLI reproduces the TB run end-to-end; HANA spatial strict mode green; no open SEV1/SEV2. \\
GA (limited) & Two consumer projects (TB plus Arabic AP). & M07, M08 closed; S01 closed; consumer-integration runbook (S03) exercised end-to-end; kill-switch green. \\
Full rollout & All three current consumers (TB, AP, Regulations). & S02 closed; two quarterly reviews sustained; business-case refresh complete; multi-consumer regression suite green. \\
\bottomrule
\end{tabularx}
\end{table}

\section{Kill-switch conditions}\label{sec:sim-kill-switch}

The kill-switch is a pre-committed rulebook: a tripped condition enters
the action column without further debate. GFAIC and sponsor are
notified within 24 hours. The conditions are evaluated continuously
(on every cycle) and do not carry dates.

\begin{table}[htbp]
\centering
\caption{Simula kill-switch.}
\label{tab:sim-kill-switch}
\footnotesize
\begin{tabularx}{\linewidth}{@{}l l X@{}}
\toprule
\textbf{Condition} & \textbf{Evaluation} & \textbf{Action} \\ \midrule
Reference run non-reproducible across two consecutive seeds (output diff $>$\,0) & per release & \textbf{Halt release}; investigate non-determinism; do not ship v1.2 until reproducibility is restored. \\
Schema-registry validation fails on $>$\,0.1\,\% of generated examples & per run & \textbf{Halt consumer delivery}; fix generation engine; re-run. \\
HANA spatial strict mode detects unknown SRID in production path & continuous & \textbf{Quarantine run}; incident-review before releasing corpus to consumer project. \\
Any consumer project (TB, AP, Regulations) declares red on a Simula-supplied corpus & continuous & \textbf{Pause Should items} S01--S03; remediate before next consumer wave. \\
Shared-platform dependency in \cref{tab:sim-shared-deps} declared red by its owner project & continuous & \textbf{Pause downstream Must}; re-plan at next steering. \\
\bottomrule
\end{tabularx}
\end{table}

\section{Shared-platform dependencies}\label{sec:sim-shared-deps}

This project depends on three shared-platform components. Each is
tracked here only as a dependency, not a commitment this project can
unilaterally adjust. Sequence is expressed in terms of the Must item
that first consumes the dependency, not in calendar time.

\begin{table}[htbp]
\centering
\caption{Shared-platform dependencies.}
\label{tab:sim-shared-deps}
\footnotesize
\begin{tabularx}{\linewidth}{@{}l l l X@{}}
\toprule
\textbf{Component} & \textbf{Owner} & \textbf{First consumed by} & \textbf{Impact if late} \\ \midrule
Unified schema registry (\path{docs/schema/}) & Shared infra & M03 & Already \texttt{migration.status=complete}; no remaining risk. \\
vLLM TurboQuant hosting (\texttt{vllm-only} policy) & Shared platform & M03 & Graceful-degradation ladder defined in the Regulations spec provides fallback; batch-generation slip tolerable. \\
HANA Cloud spatial extensions & Shared infra & M05 & No substitute; slip blocks spatial-policy certification and downstream consumer runs. \\
\bottomrule
\end{tabularx}
\end{table}

\noindent
Shared-platform dependencies are surfaced at monthly steering. If any
upstream owner declares \textbf{amber} or \textbf{red} on a component
in \cref{tab:sim-shared-deps}, the affected downstream Must milestone
is re-planned at the following steering.

\section{v1.3 expansion trigger}\label{sec:sim-v13-trigger}

The Could item C01 (SIM-V13-EXPLORE) is unlocked only when all of the
following are observed after two consecutive consumer corpus
generations under v1.2:

\begin{itemize}
\item aggregate consumer request volume for corpora outside the
current taxonomy exceeds 20\,\% of total Simula runs;
\item reproducibility and schema-validation rates sustained at or
above release thresholds;
\item business-case refresh (M08) explicitly budgets the v1.3
programme.
\end{itemize}

\noindent
If any of the three conditions fails, C01 remains in the Could column
and is re-evaluated at the next business-case refresh.

\section{Governance and business-case cadence}\label{sec:sim-governance}

\subsubsection{Project governance.}
The Simula project reports monthly to a project steering (Studio AI
Lead, Training Data PM, Pipeline Engineer) and quarterly to the
cross-project GFAIC. Monthly steering accepts or rejects amber/red
signals on Must items; GFAIC reviews compliance posture and any
cross-project dependencies (especially consumer-project risk).

\subsubsection{Business-case refresh (M08).}
The business case that authorises this project is refreshed on the
following rhythm:

\begin{itemize}
\item \textbf{Mandatory refresh}: at the first cycle after release,
and annually thereafter (ticket SIM-BC-REFRESH, priority M in
\cref{tab:sim-milestones}). Refresh packet includes consumer
demand assumptions, unit economics,
GPU cost realised vs.\ forecast, and a deviation statement
against the baseline.
\item \textbf{Trigger-based refresh}: any of
\begin{itemize}
\item consumer request volume changes by $\ge\,20$\,\% from
forecast (either direction);
\item per-run generation cost changes by $\ge\,30$\,\% (GPU
price, vLLM licensing, HANA I/O);
\item a consumer project's regulatory obligation expands
Simula's in-scope taxonomy;
\item any kill-switch in \cref{tab:sim-kill-switch} is
tripped.
\end{itemize}
\item \textbf{Owner}: Studio AI Lead. Approval routed through the
project steering and filed at
\path{docs/simula/business-case/refresh-history.yaml}.
\end{itemize}
\chapter{Agent Implementation Instructions}
\label{ch:agent-impl}

\section{Purpose and audience}
This chapter is the self-contained, intent-based instruction set for
implementing the \emph{Simula training data framework} domain inside
the unified SAP OSS platform. It is written for LLM-driven and human
implementation agents who will extend existing services to realise
the specification. It assumes only this PDF; every platform
decision, integration contract, and governance obligation the Simula
workstream depends on is restated here so that no companion document
is required. Existing sections~\ref{chap:overview}--\ref{sec:sim-roadmap}
remain the normative source for the Simula domain itself; this
chapter tells future agents \emph{how} to build what those sections
describe, \emph{where} to put it in the codebase, and \emph{how} it
plugs into the wider platform.

% !TEX root = ../master.tex
% =============================================================================
% MASTER PLATFORM CORE: Shared Intent-Based Instructions (Union)
% World-Class Version: HITL Gates, Process Integrity, and Sequence Logic
% =============================================================================

\section{Rules of Engagement for AI Agents}
\label{sec:agent-rules}

To ensure safety and architectural integrity, all implementation agents must adhere to the following Prime Directives:

\begin{enumerate}[label=\textsc{rule}-\arabic*:, leftmargin=2cm]
\item \textbf{Credential Protection.} Never log, print, or commit secrets. Use the \texttt{.secrets.example} pattern for all new integrations.
\item \textbf{Context Efficiency.} Prioritize surgical edits. Request only the minimum required file ranges to satisfy the task intent.
\item \textbf{Test-Driven Implementation.} No code is complete without a corresponding fixture in \texttt{tests/fixtures/} and a passing validation run.
\item \textbf{Self-Correction.} If a schema validation fails, the agent must diagnose the root cause (hallucination vs. schema drift) before re-attempting implementation.
\end{enumerate}

\section{Platform Overview --- One Codebase, One Instance}
\label{sec:agent-platform-overview}

\subsection{Integration Stance}
The SAP OSS platform is delivered as a single codebase that runs as a single logical instance. All cross-cutting concerns are hosted by existing services.

\begin{adrbox}{{One Codebase, One Instance}}
To minimize integration latency and eliminate distributed-system failure modes (e.g., partial partitions), we mandate a single logical instance. All cross-domain orchestration happens via the MCP Gateway layer.
\end{adrbox}

\subsection{Platform Integration Map}
\label{sec:platform-map}

\begin{figure}[htbp]
\centering
\begin{tikzpicture}[
node distance=2.5cm,
auto,
block/.style={rectangle, draw, rounded corners, fill=sapblue!10, text width=6em, text centered, minimum height=3em, font=\scriptsize},
regs/.style={rectangle, draw, rounded corners, fill=sapgrey!20, text width=7em, text centered, minimum height=4em, thick, dotted, font=\scriptsize}
]
\node (arabic) [block] {Arabic AP};
\node (tb) [block, below of=arabic] {Trial Balance};
\node (simula) [block, right of=arabic, node distance=5.5cm] {Simula Training};
\node (regs) [regs, below of=simula] {Regulations (Runtime Wrapper)};

\draw [->, thick] (arabic) -- node [align=left, font=\tiny] {Posted\\Invoices} (tb);
\draw [->, thick, dashed] (simula) -- node [align=left, font=\tiny, near start] {Training\\Examples} (arabic);
\draw [->, thick, dashed] (simula) -- node [align=left, font=\tiny, near start] {Training\\Examples} (tb);
\draw [->, thick] (regs) -| node [font=\tiny, near end, xshift=-1cm] {Governance Gate} (arabic);
\draw [->, thick] (regs) -| node [font=\tiny, near end, xshift=-1cm] {Audit \& Identity} (tb);
\end{tikzpicture}
\caption{Cross-domain interaction and unified data flow map. Detailed governance logic is defined in \cref{reg:ch:mgf}.}
\end{figure}

\section{Transactional Sequence Diagrams}
\label{sec:agent-sequences}

Implementation agents must adhere to the temporal logic defined below. Every cross-domain request MUST be brokered by the Regulations Wrapper.

\subsection{Universal Governance Handshake}
This sequence applies to \textbf{all} platform service calls.

\begin{figure}[htbp]
\centering
\begin{tikzpicture}[x=1cm, y=-0.8cm, font=\scriptsize]
% Lifelines
\draw[thick] (0,0) -- (0,7) node[below] {Persona (UI)};
\draw[thick] (3,0) -- (3,7) node[below] {MCP Gateway};
\draw[thick] (6,0) -- (6,7) node[below] {Regs Wrapper};
\draw[thick] (9,0) -- (9,7) node[below] {Domain Service};

% Interaction
\draw[->, >=latex] (0,1) -- node[above] {1. Request + Identity} (3,1);
\draw[->, >=latex] (3,2) -- node[above] {2. Validate (Pillar 1)} (6,2);
\draw[<-, >=latex] (3,3) -- node[above] {3. Auth/Allow Decision} (6,3);
\draw[->, >=latex] (3,4) -- node[above] {4. Execute Tool} (9,4);
\draw[->, >=latex] (9,5) -- node[above] {5. Log Evidence (Pillar 2)} (6,5);
\draw[<-, >=latex] (0,6) -- node[above] {6. Response + Audit ID} (3,6);

\node[draw, fill=yellow!10, text width=2.5cm, align=center] at (4.5,0.2) {MAS MGF Gate};
\end{tikzpicture}
\caption{Standard platform handshake for governed tool execution. Audit schemas are detailed in \cref{reg:ch:identity}.}
\end{figure}

\subsection{Arabic AP Golden Path}
Temporal flow for invoice processing with Arabic context.

\begin{figure}[htbp]
\centering
\begin{tikzpicture}[x=1cm, y=-0.6cm, font=\tiny]
\draw[thick] (0,0) -- (0,8) node[below] {Intake};
\draw[thick] (2.5,0) -- (2.5,8) node[below] {OCR};
\draw[thick] (5,0) -- (5,8) node[below] {VAT Engine};
\draw[thick] (7.5,0) -- (7.5,8) node[below] {Workflow};
\draw[thick] (10,0) -- (10,8) node[below] {Human (HITL)};

\draw[->] (0,1) -- node[above] {Extract} (2.5,1);
\draw[->] (2.5,2) -- node[above] {Translate} (5,2);
\draw[->] (5,3) -- node[above] {Evaluate} (5,4); % Internal loop
\draw[->] (5,5) -- node[above] {\textsc{Hold} Notification} (10,5);
\draw[<-] (5,6) -- node[above] {Manual Release} (10,6);
\draw[->] (5,7) -- node[above] {Final Posting} (7.5,7);
\end{tikzpicture}
\caption{Arabic AP sequence: Scanned PDF to Release. VAT rules are specified in \cref{ar:ch:vat}.}
\end{figure}

\section{Compute and Resource Alignment}
\label{sec:hardware-profiles}

AI algorithms are calibrated for the reference infrastructure established in the domain chapters:

\begin{table}[htbp]
\centering
\footnotesize
\caption{Infrastructure Reference Profiles.}
\label{tab:hardware-profiles}
\begin{tabularx}{\textwidth}{@{}l X X@{}}
\toprule
\textbf{Workload} & \textbf{Reference Hardware} & \textbf{Performance Target} \\
\midrule
LLM Inference & NVIDIA L40S (48GB) & AWQ 4-bit / p95 < 2s \\
Embeddings & NVIDIA A100 (40/80GB) & FP16 / p99 < 500ms \\
Simula Training & NVIDIA A100 SXM4 cluster & BF16 throughput parity \\
\bottomrule
\end{tabularx}
\end{table}

\section{Human-in-the-Loop (HITL) Critical Gates}
\label{sec:hitl-gates}

The following "Four-Eyes" gates are mandatory. AI agents must implement the "Awaiting Sign-off" state for these actions:
\begin{itemize}
\item \textbf{VAT Release:} No invoice marked as \textsc{hold} can be released without human review (see \cref{ar:ch:ap}).
\item \textbf{Anomaly Override:} Statistical outliers ($3\sigma$) cannot be dismissed by AI; only "Confirmed" or "Triaged" by human (see \cref{tb:ch:controls}).
\item \textbf{Journal Posting:} Agents may \emph{draft} journal entries, but only the Financial Controller persona can \emph{commit} to S/4HANA (see \cref{tb:ch:controls-engine}).
\end{itemize}

\section{Registry Integrity \& Live Status}
\label{sec:registry-integrity}

Before initiating any schema-dependent task, implementation agents must verify the live status of the Unified Schema Registry (see \cref{sim:chap:schema}).
\begin{adrbox}{Registry First}
If \texttt{docs/schema/registry.json} is out of sync with the codebase, the agent must halt and emit a \textsc{schema-mismatch} event.
\end{adrbox}

\section{Master Implementation Manifest (Union)}
Implementation runs in three waves across the entire platform.

\subsection*{Wave 1 --- Foundations}
\begin{itemize}
\item \textbf{Simula:} Implement HANA schema extraction pipeline (MCP discovery). See \cref{sim:chap:extraction}.
\item \textbf{Regulations:} Implement MAS MGF Pillar 1 logic (Tool allow-list). See \cref{reg:ch:mgf}.
\item \textbf{Arabic:} Implement Invoice Intake \& OCR Gateway. See \cref{ar:ch:ai}.
\item \textbf{Trial Balance:} Implement identity-attribution wiring. See \cref{reg:ch:identity}.
\end{itemize}

\subsection*{Wave 2 --- Core Domain Build-out}
\begin{itemize}
\item \textbf{Trial Balance:} Implement Controls \& Commentary Engine (3$\sigma$ anomaly logic). See \cref{tb:ch:ai,tb:ch:controls-engine}.
\item \textbf{Arabic:} Implement multi-country VAT Engine (Saudi, UAE, Bahrain, Oman). See \cref{ar:ch:vat,ar:ch:vat-impl}.
\item \textbf{Simula:} Implement Taxonomy Generation \& Agentic Data Synthesis. See \cref{sim:chap:taxonomy,sim:chap:generator}.
\item \textbf{Trial Balance:} Realise 28-state BPMN workflow (\texttt{tb-review-glo}). See \cref{tb:ch:workflow}.
\end{itemize}

\subsection*{Wave 3 --- Hardening}
\begin{itemize}
\item \textbf{Regulations:} Implement Per-Request Identity Attribution. See \cref{reg:ch:identity}.
\item \textbf{Trial Balance:} Deliver Human Review Dashboard. See \cref{tb:ch:review-dashboard}.
\item \textbf{Simula:} Deliver test fixtures and testing at 85\% coverage. See \cref{sim:ch:test-fixtures,sim:ch:testing}.
\item \textbf{Regulations:} Wire Emergent-Capability monitoring. See \cref{reg:ch:monitoring}.
\end{itemize}

\section{Governance Obligations (MAS MGF)}
The regulations domain wraps all service calls at runtime. Agents must implement the following hooks in every domain:
\begin{itemize}
\item \textbf{Pillar 1:} Tool allow-list enforcement (deny-by-default).
\item \textbf{Pillar 2:} Traceability via the identity attribution envelope.
\item \textbf{Pillar 3:} Metric emission to the monitoring sink.
\end{itemize}
Detailed pillar requirements are defined in \cref{reg:ch:mgf}.

\section{FMEA: Failure Mode and Effects Analysis for AI}
\label{sec:fmea-ai}

The following analysis identifies critical AI failure modes within the Simula pipeline and their deterministic mitigations.

\begin{table}[htbp]
\centering
\footnotesize
\caption{FMEA for Simula AI Pipeline.}
\label{tab:simula-fmea}
\begin{tabularx}{\textwidth}{@{}l X X X@{}}
\toprule
\textbf{Failure Mode} & \textbf{Symptom} & \textbf{Impact} & \textbf{Technical Mitigation} \\
\midrule
Critic Collusion & Both critics pass a low-quality SQL query. & High (Model Bias) & Inject a 'Gold Standard' probe query every 100 samples. \\
Schema Contraction & LLM ignores complex JOIN paths in the taxonomy. & Medium (Coverage) & Taxonomic coverage check against Graph-DB baseline. \\
Elo-Score Drift & Calibrated complexity scores diverge from human ratings. & Medium (UI/UX) & Monthly Elo-recalibration against human-label set. \\
\bottomrule
\end{tabularx}
\end{table}

\section{This domain in context}
\label{sec:agent-domain-context}
Simula is the training-data workstream among the four domains. Its
primary persona is the Training Engineer, whose working view is the
Simula console surfaced through the shared workspace shell and gated
by Simula entitlements. The domain's entry points are the schema
sets and fixtures exported by Arabic AP, TB, and regulations, plus
the HANA Cloud catalogue read via the extraction pipeline; its exit
points are \texttt{TrainingExample} records emitted to the training
corpus and Elo-calibrated complexity reports. Every artefact is
validated against the Simula-owned schemas in
\texttt{structured/schema/} before it leaves the pipeline. The
domain is read-only with respect to the other three domains at the
schema layer: it never mutates their records, and it never calls
their runtime MCP tools. The horizontal regulations wrapper governs
Simula calls like any other domain.

\section{Schema contribution}
\label{sec:agent-schema-contrib}
Simula owns the schemas catalogued in
chapter~\ref{chap:schema}, including the training-example record,
the schema-registry record, the taxonomy structure, and the
enumerations listed there. These schemas live under
\texttt{docs/latex/specs/simula/} as normative text and under
\texttt{structured/schema/} as the JSON Schema artefacts consumed by
the pipeline at \texttt{src/training/pipeline/}. Registration in the
unified registry is additive: each schema is listed in
\texttt{src/intelligence/ai-core-pal/data\_products/registry.yaml}
with its URI, owning domain (\texttt{simula}), security class
(\texttt{confidential}), and LLM routing (\texttt{vllm-only}),
following the ODPS~4.1 pattern already documented for HANA training
data products in that file. Simula also consumes the schemas owned
by Arabic, TB, and regulations as inputs; the contract is strictly
schema-in, \texttt{TrainingExample}-out. Simula performs no schema
merges; the two \texttt{requirement.schema.json} files in TB and
regulations remain distinct when Simula reads them.

\section{Persona and MCP allow-list}
\label{sec:agent-persona}
\begin{itemize}
\item \textbf{Persona id:} \texttt{training-engineer}.
\item \textbf{Roles:} schema extraction, taxonomy authoring,
training-data synthesis supervision, critic-calibration review,
corpus export.
\item \textbf{Entitlements:} read/write on the training-example
record, schema-registry record, and taxonomy structure; read on
every other domain's schemas (Arabic, TB, regulations);
run-execute on Simula pipeline stages described in
chapter~\ref{chap:generator}.
\item \textbf{UI route:} \texttt{/simula/\{tenant\}/console} on the
shared workspace shell; surfaced only to callers whose
entitlement set includes \texttt{training-engineer}.
\item \textbf{MCP tool allow-list:}
\texttt{simula.schema.extract}, \texttt{simula.taxonomy.build},
\texttt{simula.generate.examples},
\texttt{simula.critic.calibrate},
\texttt{simula.evaluate.coverage},
\texttt{simula.evaluate.diversity},
\texttt{simula.corpus.export}. All other MCP tools are denied by
default per MAS MGF pillar~1.
\end{itemize}

\section{Implementation tasks for this domain}
\label{sec:agent-tasks}
The tasks below are intent-based: each task states the outcome a
future agent must achieve and the location in the existing codebase
where the work belongs. None of them contains code; each task binds
to the Simula sections that define the required behaviour.

\subsubsection{Task S1 --- HANA schema extraction.}
\textbf{Intent.} Realise the HANA schema extraction pipeline in
chapter~\ref{chap:extraction}, covering MCP discovery, table
enumeration, column extraction, and registry population.

\begin{adrbox}{Why Direct SQL Metadata?}
We chose direct SQL-metadata extraction over high-level MCP discovery for the baseline to ensure we capture SAP-specific table types (Global Temporary, Row-Store) that standard MCP tools often omit.
\end{adrbox}

\textbf{In-scope.} Each of the four stages in the pipeline figure
(\ref{fig:extraction-pipeline}); schema registry writer; per-table
dimension/measure classification.
\textbf{Out-of-scope.} Non-HANA backends.
\textbf{Paths to extend.} \texttt{src/training/pipeline/} extraction
module; registry writer that emits entries conforming to the
schema-registry record.
\textbf{Definition of done.} Running extraction against a HANA
Cloud catalogue produces a registry populated with every in-scope
table; registry entries validate against the schema-registry record.

\textbf{Autonomous Verification (Agent Logic):}
\begin{lstlisting}[language=jsonx]
// Verify Task S1
{
"test_fixture": "tests/fixtures/simula/hana_mock.sql",
"expected_outcome": "table_count == 14",
"validation_command": "python src/simula/extract.py --verify-registry"
}
\end{lstlisting}

\textbf{Verification.} Pipeline dry-run against the fixture store
described in chapter~\ref{ch:test-fixtures} produces the expected
registry entries.

\subsubsection{Task S2 --- Taxonomy generation (Algorithm 1).}
\textbf{Intent.} Implement Algorithm~1 (reasoning-driven taxonomy
generation) from chapter~\ref{chap:taxonomy} in full, including all
four steps (factor identification, breadth-first expansion,
best-of-N proposal, critic refinement).
\textbf{In-scope.} Per-step module with the checkpoints
in~\ref{fig:taxonomy-stages}; taxonomy structure writer.
\textbf{Out-of-scope.} Manual taxonomy authoring UI.
\textbf{Paths to extend.} Taxonomy engine colocated with the
pipeline.
\textbf{Definition of done.} Running Algorithm~1 against an extracted
registry produces a taxonomy that validates against the taxonomy
structure schema and passes the coverage check in
chapter~\ref{chap:evaluation}.
\textbf{Verification.} Taxonomy engine unit tests pass; the
coverage evaluator reports coverage within the chapter's target
band.

\subsubsection{Task S3 --- Agentic data synthesis (Algorithm 2).}
\textbf{Intent.} Implement Algorithm~2 (agentic data synthesis) from
chapter~\ref{chap:generator}, including all four steps (mix
sampling, meta-prompt generation, complexification, example
generation) and the double-critic validation.
\textbf{In-scope.} The agent roster implied by
\ref{fig:generator-stages}; vLLM TurboQuant routing for all LLM
calls; complexification step and the two critic agents.
\textbf{Out-of-scope.} Non-agentic synthesis paths.
\textbf{Paths to extend.} Generation engine in
\texttt{src/training/pipeline/}; vLLM client wrapper.
\textbf{Definition of done.} Running Algorithm~2 against a taxonomy
produces \texttt{TrainingExample} records that pass both critics and
validate against the training-example record schema.
\textbf{Verification.} Generation unit tests pass on fixture inputs;
double-critic agreement rate meets the threshold in
chapter~\ref{chap:evaluation}.

\subsubsection{Task S4 --- Environment and CLI.}
\textbf{Intent.} Deliver the environment and CLI surfaces specified
in chapter~\ref{chap:config}, including the configuration table, CLI
usage, and output files listed in~\ref{tab:outputs}.
\textbf{In-scope.} Configuration loader; CLI subcommands for each
pipeline stage; output-file writers.
\textbf{Out-of-scope.} IDE integrations.
\textbf{Paths to extend.} CLI module in
\texttt{src/training/pipeline/}.
\textbf{Definition of done.} Each configuration key in the
configuration table is read and respected; the CLI runs end-to-end
and produces every output file in~\ref{tab:outputs}.
\textbf{Verification.} CLI smoke tests pass; output files match the
schemas declared for them.

\subsubsection{Task S5 --- Evaluation framework.}
\textbf{Intent.} Deliver the evaluation framework in
chapter~\ref{chap:evaluation}, including calibrated complexity
scoring (Algorithm~3, Elo), taxonomic coverage evaluation,
embedding-based diversity, and critic-calibration validation.
\textbf{In-scope.} Each evaluator as a runnable function; report
writers.
\textbf{Out-of-scope.} Real-time evaluation dashboards beyond what
the chapter defines.
\textbf{Paths to extend.} Evaluation module in
\texttt{src/training/pipeline/}.
\textbf{Definition of done.} Each evaluator runs against a fixture
corpus and emits a report that validates against the corresponding
report schema.
\textbf{Verification.} Evaluator unit tests pass; reports reproduce
under fixed seeds.

\subsubsection{Task S6 --- Test fixtures and testing.}
\textbf{Intent.} Deliver the test fixtures specification in
chapter~\ref{ch:test-fixtures} and the testing specification in
chapter~\ref{ch:testing} at the coverage thresholds
in~\ref{tab:simula-coverage}.
\textbf{In-scope.} Invoice, OCR, taxonomy, and training-example
fixtures; unit tests for taxonomy traversal and generation; CI
wiring.
\textbf{Out-of-scope.} Performance benchmarking infrastructure.
\textbf{Paths to extend.} Test directory colocated with the Simula
pipeline code.
\textbf{Definition of done.} Coverage meets~\ref{tab:simula-coverage};
CI gates block regressions.
\textbf{Verification.} \texttt{make spec-simula} remains green; the
Simula test entry points all run.


\section{Governance obligations applied to this domain}
\label{sec:agent-governance}
The regulations domain wraps Simula at runtime. Four obligations
apply without exception; each is stated as a concrete hook in the
Simula codebase.

\begin{itemize}
\item \textbf{MAS MGF four pillars} (see regulations
chapter~3 in
\texttt{docs/latex/specs/regulations/chapters/03-mgf-framework.tex}).
Pillar~1 (assess and bound): every Simula MCP tool is declared
in the allow-list in \S\ref{sec:agent-persona}, deny-by-default.
Pillar~2 (design and deploy): the pipeline stages in
chapter~\ref{chap:generator} are traceable end-to-end via the
per-request identity attribution described below. Pillar~3
(operate and monitor) and pillar~4 (govern and disclose) are
satisfied by emitting the evidence records produced by the
Simula evaluation framework (chapter~\ref{chap:evaluation}) into
the shared governance store.
\item \textbf{Identity attribution} (regulations chapter~10,
\texttt{10-identity-attribution.tex}). Every MCP call and every
pipeline run carries the persona id, agent id, request id, and
tenant id. The gateway is the attachment point; Simula code
propagates the identity envelope without mutating it.
\item \textbf{Emergent-capability monitoring} (regulations
chapter~8, \texttt{08-monitoring-spec.tex}). Extraction,
taxonomy generation, synthesis, and evaluation steps are
instrumented with the monitoring metrics specified in that
chapter; metrics flow to the monitoring sink defined there.
\item \textbf{Conformance tooling} (regulations chapter~6,
\texttt{06-conformance-tooling.tex}). The Simula service
exposes AI Verify 2.0 plug-in hooks and Moonshot CI/CD 1.1.0
benchmark entry points as required by that chapter.
\end{itemize}

\section{Cross-spec integration --- full detail}
\label{sec:agent-integration}

\subsection{With Arabic AP}
Simula consumes Arabic's schemas and fixtures as training input. The
contract is schema-in, \texttt{TrainingExample}-out: Simula reads
\texttt{invoice.schema.json}, \texttt{vat-checklist.schema.json}, and
the Arabic fixture store; it produces \texttt{TrainingExample}
records. Simula does not call Arabic MCP tools at runtime and does
not share state with the Arabic workflow. No cross-domain
operational envelope is introduced.

\subsection{With Trial Balance}
Simula consumes TB's schemas and fixtures as training input. The
contract is schema-in, \texttt{TrainingExample}-out: Simula reads
the TB schema set and the TB fixture store; it produces
\texttt{TrainingExample} records. Simula does not call TB MCP tools
at runtime and does not share state with the TB workflow.

\subsection{With Regulations}
Regulations wraps Simula at runtime. Every Simula MCP call is
brokered through the regulations governance gate, which applies
pillars~1--4 of MAS MGF, attaches the identity envelope, and routes
emergent-capability metrics to the monitoring store. The Simula
service code does not call regulations directly; the gateway layer
enforces this. Simula also consumes regulations' schemas
(requirement, conformance-tool, monitoring metrics, identity and
audit) as training input.

\subsection{With Simula}
Simula exports the training-example record, the schema-registry
record, the taxonomy structure, and the enum vocabularies under
\texttt{structured/schema/} (chapter~\ref{chap:schema}); registers
the MCP tools enumerated in \S\ref{sec:agent-persona}; and applies
the four governance hooks in \S\ref{sec:agent-governance} to every
call it serves. External consumers see Simula only through these
contracts and through the corpus it emits; there are no private
back-channels.

\section{Wave assignment and sequencing}
\label{sec:agent-waves}
Implementation runs in three waves.

\begin{itemize}
\item \textbf{Wave 1 --- foundations.} Unified schema registry
entry; persona-gated surface for the Training Engineer;
identity attribution envelope flowing through the gateway; HANA
extraction pipeline (Task S1). Prerequisites: the regulations
chapter~10 identity schema must be frozen.
\item \textbf{Wave 2 --- Simula core build-out.} Tasks S2, S3, S4,
S5. Prerequisites: Wave~1 complete; vLLM TurboQuant endpoint
reachable; Arabic, TB, and regulations schema sets published.
\item \textbf{Wave 3 --- hardening and conformance.} Task S6 and
full conformance-tooling hook-up. Prerequisites: Wave~2
complete; regulations Wave~2 conformance tooling integrated;
monitoring sink live.
\end{itemize}

Simula occupies Wave~1 for its extraction pipeline (because other
domains depend on Simula's fixtures being reachable), Wave~2 for its
core build-out, and Wave~3 for its testing and conformance surfaces.

\section{Acceptance criteria (platform-wide)}
\label{sec:agent-acceptance}
When all four chapter~12 instruction sets are implemented, the
platform satisfies the following invariants:

\begin{enumerate}
\item \textbf{One instance.} A single codebase runs a single
logical instance; no domain ships as a separate product.
\item \textbf{Persona-gated surfaces.} The four personas each
reach a tailored UI on the shared shell; no route, MCP tool, or
data product is visible outside its persona's entitlement set.
\item \textbf{Unified registry.} Every schema in use is listed in
the ODPS~4.1 registry under its owning domain; validators
resolve \texttt{\$ref}s through the unified registry only.
\item \textbf{No schema merges.} The two
\texttt{requirement.schema.json} files in TB and regulations
remain distinct; no cross-domain composite schemas are produced.
\item \textbf{No cross-domain operational envelope.} Runtime
orchestration stays within each domain's service; only Simula
consumes other domains' outputs, and only as training input.
\item \textbf{Governance wraps all three non-regulations domains.}
Arabic, TB, and Simula service calls pass through the
regulations runtime wrapper without exception.
\end{enumerate}

\section{References}
\label{sec:agent-references}
Within-spec references:
chapter~\ref{chap:overview} (overview),
chapter~\ref{chap:schema} (data schema),
chapter~\ref{chap:extraction} (HANA extraction pipeline),
chapter~\ref{chap:taxonomy} (taxonomy engine, Algorithm~1),
chapter~\ref{chap:generator} (generation engine, Algorithm~2),
chapter~\ref{chap:config} (environment and CLI),
chapter~\ref{chap:evaluation} (evaluation framework),
chapter~\ref{ch:test-fixtures} (test fixtures),
chapter~\ref{ch:testing} (testing specification),
chapter~\ref{sec:sim-roadmap} (project roadmap).

Peer chapter~12 instruction sets for the other three domains (each
is self-contained; consult directly when operating in that domain):
\begin{itemize}
\item Arabic AP:
\texttt{docs/latex/specs/arabic/chapters/12-implementation-instructions.tex}.
\item Trial Balance:
\texttt{docs/latex/specs/tb/chapters/12-implementation-instructions.tex}.
\item Regulations:
\texttt{docs/latex/specs/regulations/chapters/12-implementation-instructions.tex}.
\end{itemize}

Shared platform artefacts:
\begin{itemize}
\item ODPS~4.1 registry:
\texttt{src/intelligence/ai-core-pal/data\_products/registry.yaml}.
\item Schema pipeline and resolver:
\texttt{src/intelligence/schema\_pipeline/}.
\item Training pipeline (Simula):
\texttt{src/training/pipeline/}.
\item Shared workspace shell:
\texttt{generativeUI/training-webcomponents-ngx}.
\end{itemize}

% --- Incorporated Universal Chapters ---
\chapter{Global Wave-Dependency Manifest}
\label{ch:global-dependencies}

\section{Critical Path Visualization (Waves 1--3)}
The following chart defines the normative sequencing for implementation agents. A task in a later wave cannot be marked \textsc{complete} until its prerequisites in an earlier wave are validated.

\begin{figure}[htbp]
\centering
\begin{tikzpicture}[
node distance=1.5cm,
auto,
wavebox/.style={rectangle, draw, rounded corners, fill=sapmidnight!10, text width=10em, text centered, minimum height=3em, font=\scriptsize\bfseries},
dep/.style={->, >=latex, thick, sapblue},
critical/.style={->, >=latex, very thick, sapamber}
]
% --- Wave 1: Foundations ---
\node (r1) [wavebox, fill=sapgrey!20] {\textcolor{sapmidnight}{REG-W1: Identity \& Audit Schema}};
\node (s1) [wavebox, below of=r1, yshift=-0.5cm] {SIM-W1: HANA Extraction Pipeline};
\node (a1) [wavebox, below of=s1, yshift=-0.5cm] {ARA-W1: Invoice Intake \& OCR};
\node (t1) [wavebox, below of=a1, yshift=-0.5cm] {TB-W1: Identity Wiring};

% --- Wave 2: Core Build-out ---
\node (r2) [wavebox, right of=r1, xshift=4.5cm, fill=sapmidnight!20] {\textcolor{sapmidnight}{REG-W2: Governance Runtime}};
\node (s2) [wavebox, right of=s1, xshift=4.5cm] {SIM-W2: Taxonomy \& Synthesis};
\node (a2) [wavebox, right of=a1, xshift=4.5cm] {ARA-W2: Multi-Country VAT};
\node (t2) [wavebox, right of=t1, xshift=4.5cm] {TB-W2: Controls \& Workflow};

% --- Wave 3: Hardening ---
\node (r3) [wavebox, right of=r2, xshift=4.5cm, fill=sapmidnight!30] {\textcolor{sapmidnight}{REG-W3: Conformance Suite}};
\node (s3) [wavebox, right of=s2, xshift=4.5cm] {SIM-W3: 85\% Coverage Test};
\node (a3) [wavebox, right of=a2, xshift=4.5cm] {ARA-W3: Calibration Pilot};
\node (t3) [wavebox, right of=t2, xshift=4.5cm] {TB-W3: HITL Dashboard};

% --- Cross-Domain Dependencies (The Critical Path) ---
\draw [critical] (r1) -- node [above, font=\tiny, align=center] {Identity Envelope\\Invariants} (a1.north east);
\draw [critical] (r1) -- (t1.north east);
\draw [dep] (s1.east) -- node [above, font=\tiny, xshift=-0.5cm] {Training Fixtures} (t2.west);
\draw [dep] (s1.east) -- (a2.west);
\draw [dep] (r2.south) -- node [left, font=\tiny, yshift=0.3cm] {Policy Gate} (a2.north);
\draw [dep] (t2.east) -- node [above, font=\tiny] {Audit Evidence} (r3.west);

% --- Wave Transitions ---
\draw [->, dashed, sapgrey] (r1.east) -- (r2.west);
\draw [->, dashed, sapgrey] (r2.east) -- (r3.west);

% --- Legend ---
\node[draw, sapamber, very thick, label=right:\scriptsize Critical Prerequisite] at (0,-7.5) {};
\node[draw, sapblue, thick, label=right:\scriptsize Data Dependency, xshift=4cm] at (0,-7.5) {};

\end{tikzpicture}
\caption{Cross-specification dependency graph: The blueprint for platform rollout.}
\end{figure}

\section{Inter-Domain Prerequisite Matrix}

\begin{table}[htbp]
\centering
\footnotesize
\caption{Normative Cross-Domain Prerequisites.}
\label{tab:prerequisites}
\begin{tabularx}{\textwidth}{@{}l l l X@{}}
\toprule
\textbf{Target Task} & \textbf{Prerequisite} & \textbf{Ref ID} & \textbf{Rationale} \\
\midrule
\rowcolor{sapmidnight!5}
\textbf{ARA-W1} & REG-W1 (Identity) & \texttt{REG-IDENTITY-01} & Invoice intake must be audit-traced from Day 1. \\
\textbf{TB-W2}  & SIM-W1 (Extraction) & \texttt{SIM-EXTRACT-01} & Anomaly detection requires HANA schema metadata. \\
\rowcolor{sapmidnight!5}
\textbf{REG-W2} & SIM-W1 (Registry) & \texttt{SIM-REGISTRY-01} & Governance monitoring relies on the schema registry. \\
\textbf{ARA-W2} & REG-W2 (Pillar 1) & \texttt{REG-MGF-P1} & VAT release is an HITL gate governed by MAS Pillar 1. \\
\rowcolor{sapmidnight!5}
\textbf{TB-W3}  & ARA-W2 (Posted Invoices) & \texttt{ARA-AP-POST} & Review dashboard depends on posted AP data from Arabic domain. \\
\bottomrule
\end{tabularx}
\end{table}

\begin{adrbox}{Dependency Enforcement}
Implementation agents are strictly forbidden from "Fast-Forwarding" into Wave 2 if Wave 1 critical prerequisites (marked in \textcolor{sapamber}{\textbf{Gold}}) are in a \textsc{gap} or \textsc{partial} state.
\end{adrbox}



\end{document}
